\documentclass{article}
\usepackage{hyperref}
\hypersetup{
    urlcolor=blue,
    citecolor=blue,
    colorlinks=true
%    colorlinks=true
}

%    filecolor=magenta,      
%    urlcolor=cyan,
%    pdftitle={Overleaf Example},
%    pdfpagemode=FullScreen,
%    }
\usepackage{comment}
\usepackage{tcolorbox}
\usepackage{capt-of}
\usepackage{amsmath}
\usepackage{algorithm}
\usepackage{algpseudocode}
\usepackage{cancel}
\usepackage{amsfonts} 
\usepackage{textcomp}
\usepackage{amssymb}
\usepackage{subfig}
\usepackage{caption}
\usepackage{graphicx}
\usepackage{wrapfig}
\usepackage{float}
\usepackage{cancel}
\usepackage{booktabs}
\usepackage{lipsum}
\usepackage{enumerate}
\usepackage{authblk}
\usepackage{makecell}
\usepackage[a4paper, total={6.5in, 9.5in}]{geometry}
\usepackage[utf8]{inputenc}
\usepackage[]{apacite}
\bibliographystyle{apacite}
\usepackage{lineno}
\usepackage{xcolor}
\linenumbers
\DeclareMathOperator*{\argminA}{arg\,min} 
\title{Semantic representations in episodic memory enhance recall and compositional consolidation}
\author[1]{Albert Albesa-González}
\author[1]{Claudia Clopath} 
\affil[1]{Department of Bioengineering, Imperial College London, London, UK}
\begin{document}
\maketitle
\begin{abstract}
\noindent Episodic and semantic memory are classically thought to play distinct roles: episodic memory encodes unique experiences, while semantic memory generalizes across them. Current conceptualizations of episodic and semantic memory interactions emphasize a one-way consolidation from episodic traces in the medial temporal lobe (MTL) to semantic knowledge in neocortex (CTX). However, this tradition has left largely unexplored how semantic memory may affect episodic encoding of new memories. Here we introduce a cognitive model in which the code used in episodic memories shifts from purely sensorial to including explicit semantic representations of stored events. Simultaneously, we propose a computational circuit model of how such a cognitive strategy could be implemented in the brain using biologically-plausible learning rules. We show that increased sparsity during replay enables neocortex to extract compositional structure from overlapping episodes, which creates a dictionary of inter-connected concepts in semantic memory. Furthermore, we show that spontaneous activity in neocortical areas can imprint the abstracted representations into the medial temporal lobe, giving rise to concept-like cells. This bidirectional interaction improves episodic recall and replay fidelity, and  facilitates the consolidation of higher-order representations based on previous semantic knowledge. The model accounts for behavioural advantages of schema-congruent learning, the emergence of concept neurons, and enhanced memory performance for semantically familiar stimuli. Together, our results provide a mechanistic account of how episodes and concepts reinforce each other, extending standard consolidation theories and suggesting a cooperative framework where semantic knowledge scaffolds episodic encoding, which in turn favours compositional abstraction.
\end{abstract}
\section*{Introduction}
The distinction between episodic and semantic memory has been central to cognitive neuroscience since its first formalization \shortcite{tulving_1972}. In its original view, episodic memory captures individual experiences, while semantic memory accumulates generalized knowledge that emerges from repeated exposure to similar episodes \shortcite{squire_2004}. One influential computational explanation for this division is provided by the Complementary Learning Systems (CLS) framework. This theory posits that a fast-learning system (the hippocampus and the rest of the medial temporal lobe, MTL) specializes in rapidly encoding episodic memories. In contrast, a slow-learning neocortical (CTX) system works in concert with the MTL to consolidate and generalize memories \shortcite{mcclelland_1995, oreilly_2014}.
\newline\newline
One of the central assumptions of consolidation theories is that replay events originating in the hippocampus provide a teaching signal to neocortex. This idea has been remarkably successful in explaining how semantic knowledge can gradually emerge from episodic experience and remains a cornerstone of Complementary Learning Systems theory. However, the precise computational mechanism underlying replay-driven learning remains unclear. On the one hand, several influential high-level models are strongly inspired by the machine learning paradigm \shortcite{saxe_2019, singh_2022, sun_2023, spens_2024}. As a result, they often rely on a predefined distinction between input, hidden, and output representations as the backbone of cortical learning. How such distinctions could be stored within replay events and translated into appropriate plasticity signals remains unknown. Furthermore, these approaches typically assume that hippocampal replay faithfully reinstates both the cortical inputs and target representations required for learning, an assumption that is difficult to reconcile with evidence that replay is sparse, reconstructive, and only partially reinstates previous activity \shortcite{swanson_2020, stoianov_2022}. Finally, learning commonly relies on error-driven or backpropagation-like mechanisms whose biological implementation remains debated \shortcite{lillicrap_2020}. On the other hand, more biologicall grounded models \shortcite{tome_2022, chrysanthidis_2022, guerreiro_2024} often assume a pre-existing shared code between hippocampus and cortex, sometimes implemented through direct one-to-one mappings between the two regions. Consequently, these models primarily address how associations between representations can be transferred from hippocampus to cortex, rather than how semantic representations themselves emerge from episodic experience. In addition, the complexity of the memories they can accommodate is limited, requiring highly decorrelated or orthogonal episodic patterns to ensure stable learning and recall. Together, these considerations suggest the need for an account in which replay provides candidate structure rather than predefined targets, semantic representations emerge rather than being assumed, learning relies only on local plasticity mechanisms, and replay can be generative rather than a verbatim reinstatement of previous experience.
\newline\newline
A second unresolved tension between theory and experiments concerns the nature of the episodic code itself. Classical theories of associative memory have long emphasized that storage capacity is maximized when memories are sparse, decorrelated, and approximately orthogonal to one another \shortcite{hopfield_1982, tsodyks_1988, treves_1991}. In contrast, a growing body of experimental evidence suggests that representations in the medial temporal lobe contain substantial semantic and relational structure. These include concept cells that respond selectively to abstract concepts across different sensory modalities and contexts \shortcite{quiroga_2005, quiroga_2012, courellis_2024}, as well as neurons that encode latent variables and higher-order task structure \shortcite{aronov_2017, sun_2020}. Such representations necessarily introduce structured overlap between memories, apparently violating the conditions under which episodic memory is expected to operate most efficiently. What's more, semantically meaningful material is paradoxically often remembered more accurately than arbitrary information \shortcite{murray_1976, tse_2007, lin_2021}. This raises two fundamental questions: what is the origin of these semantic representations in the episodic memory circuit, and what computational role do they serve?
\newline\newline
All in all, while views challenging a strict dichotomy between semantic and episodic memory are almost as old as the semantic-episodic distinction itself \shortcite{mckoon_1986, greenberg_2010, gentry_2024}, we still lack a mechanistic understanding of how or why these two memory systems interact. Nevertheless, several recent developments point towards a different perspective on episodic-semantic interactions. Computational work has suggested that hippocampal replay preferentially supports the learning of statistical regularities and generalizable structure across experiences \shortcite{schapiro_2017, sun_2023}. This view is consistent with an emerging interpretation of replay as a generative process that reorganizes and samples from episodic experience, rather than providing a verbatim teaching signal for cortical learning \shortcite{swanson_2020, diekmann_2023}. Under this perspective, neocortex may act less as a recipient of replayed memories and more as an interpreter of replayed structure, extracting a set of reusable representations that efficiently explain recurring replay events. Such a process bears a strong conceptual resemblance to dictionary learning frameworks, in which complex activity patterns are decomposed into sparse combinations of reusable latent features \shortcite{olshausen_1996, bricken_2023}. A complementary set of observations concerns the flow of information in the opposite direction. While consolidation theories have traditionally emphasized hippocampus-to-cortex transfer \shortcite{moscovitch_2021}, cortex-to-hippocampus teaching has recently been highlighted as an important and largely unexplored aspect of memory consolidation \shortcite{liu_2024}. Related ideas have begun to emerge in computational accounts that assign an increasingly active role to cortical semantic knowledge during replay \shortcite{van_2020} and episodic encoding \shortcite{spens_2024_b}. Furthermore, experimental work has revealed mechanisms through which cortical activity can shape medial temporal lobe representations, including cortex-initiated hippocampal replay \shortcite{rothschild_2017}, strengthening of corticofugal projections from cortical assemblies to medial temporal circuits \shortcite{lee_2023}, and neocortical signals that control the overlap and linking of episodic memories \shortcite{de_sousa_2026}. Together, these observations suggest a framework in which replay is used not simply to transfer memories between systems, but to discover reusable semantic structure from episodic experience and subsequently use that structure to shape future episodic representations.
\newline\newline
Here, we propose a model of episodic-semantic memory interactions that addresses both of the challenges outlined above. The model has a circuit implementation based on hebbian local plasticity, with neocortex as the primary substrate for semantic memory and the medial temporal lobe as the primary substrate for episodic memory. First, we put forward dictionary learning as the primary algorithm of semantic consolidation. Under this view, the role of \textit{episodic replay} is not to provide targets for cortical learning, but to generate candidate structure from which neocortex can discover recurring atomic representations, termed \textit{concepts}, that support a compositional interpretation of future experience. Second, we propose that these learned cortical representations can be transferred back into the episodic memory system through a complementary process we term \textit{semantic replay}. This mechanism allows semantic representations to become part of the episodic code itself, providing a potential explanation for concept-like and latent-variable representations observed in medial temporal lobe circuits. Then, we show that these two processes in the circuit jointly support the emergence of semantic knowledge from episodic experience and the incorporation of semantic structure into episodic memories. We further demonstrate that the resulting dynamics reproduce key experimental observations, including concept-cell-like responses, increasing overlap between semantically related memories, and recent findings in synaptic-specific changes associated with consolidation. We then investigate the computational role of semantic representations in episodic memory and find that they act as an error-correcting scaffold that improves recall, stabilizes replay dynamics, and increases the signal-to-noise ratio of replayed activity patterns. Together, these results provide a mechanistic account of how semantic knowledge can emerge from episodic experience and, in turn, shape the organization and consolidation of future memories.
\newpage
\section*{Results}
\subsection*{A model for semantic structure in episodic memory based on dictionary learning}
We introduce a cognitive model of semantic and episodic memory (Fig. \ref{fig:diagram}a). The model starts from a simple psychological observation: episodic memories often contain both sensory detail and semantic content. For example, let's consider multiple episodes of someone having their favourite food: \textit{pizza}. These may differ in appearance, smell, taste, and context, but still share semantic structure such as the presence of \textit{pizza}, \textit{cheese}, and \textit{tomato}. Our cognitive model captures this by positing that episodic traces contain two components (Fig. \ref{fig:diagram}a): a sensory component, preserving event-specific perceptual detail, and a semantic component, representing reusable concepts shared across events.
\newline\newline
At the circuit level, however, the presence of semantic representations in episodic memory raises a computational problem. The medial temporal lobe (MTL), a region crucial for episodic memory, has been shown to contain concept-like information. While classical attractor models of episodic memory achieve maximal storage capacity when patterns are maximally pattern, semantic representations introduce structured overlap between episodes. At the same time, behavioural evidence seems to point in the opposite direction of the thoeretical prediction: episodes that can be semantically comprehended are better recalled, and lesions in MTL regions like the hippocampus also impair semantic recall on already consolidated information. How is it possible that episodic memory system  can not only tolerate this structured overlap but actually seems to benefit from it? To address this question, we build a complementary circuit model to study how semantic information can coexist with episodic storage in MTL. Following standard mappings, in the model neocortex (CTX) contains semantic memory and MTL episodic memory (Fig. \ref{fig:diagram}c). MTL is divided into MTL-sensory, driven by sensory input, and MTL-semantic, driven by semantic cortical input. These populations are recurrently connected, allowing sensory and semantic representations to be bound into individual episodic traces. All regions are modeled as binary networks with homeostasis (see Technical Box \ref{box:hbn} and Methods), for their analytical tractability and direct connection with semantic learning.
\newline\newline
A related question that can be asked with the circuit model is how are the cortical semantic representations formed in the first place. A common framing is that replayed cortical reactivations are used as a teaching signal, which can then be back-propagated. However, this often leaves implicit how the reactivated episodic pattern interacts with cortical activity. Which cortical neurons provide the input? Which ones define the target? Which ones remain plastic? Even biologically motivated alternatives to backpropagation, such as Contrastive Hebbian Learning or Predictive Coding, still require a structured distinction between input, hidden, and output components during learning.
\newline\newline
Here, this problem is avoided this by proposing replay as a mechanism for dictionary learning (Technical Box \ref{box:replay}). Under this view, the role of r\textit{episodic replay} is to provide episodic patterns that CTX must interpret, which is achieved via Hebbian plasticity and competitive learning. This avoids backpropagation-like credit assignment and the need to predefine cortical targets. Furthermore, the replayed pattern is not a verbatim copy of awake experience, but a generalized episodic pattern. When projected to CTX, this pattern either refines existing representations or recruits new neurons if no previous representation is sufficiently close. 
\newline\newline
The same mechanism can be used to explain how CTX can transfer the learned representations to MTL. If recurrent MTL dynamics expose episodic structure to CTX, recurrent CTX dynamics can expose semantic structure to MTL. We call this process \textit{semantic replay} because, like episodic replay, is internally generated, begins from random activity, and recovers attractor states through recurrent dynamics. However, because these attractors are shaped by CTX recurrent connections, they reflect semantic rather than episodic structure. Projecting them to MTL-semantic provides a local learning signal through which concept-like representations can enter the episodic circuit.
\newline\newline
Together, episodic and semantic replay implement in the context of the cognitive model (Fig. \ref{fig:diagram}b,d). Episodic replay supports \textit{episodic-to-semantic consolidation}, allowing CTX to extract a cortical dictionary from episodic structure. Semantic replay supports \textit{semantic-to-episodic consolidation}, allowing MTL-semantic to incorporate the learned cortical vocabulary. We next validate this learning algorithm in simulations of the circuit model. We then use the model to ask how can semantic overlap interact in a beneficial way with episodic attractor dynamics, supporting memory and further consolidation.
\begin{figure}[h!]
    \centering
    \includegraphics[width=\linewidth]{Figures/Figure_1.png}
    \caption{\textbf{We propose a cognitive model in which semantic memory shapes representations in episodic memory, and an associated circuit implementation.} \textbf{a}: Cognitive model of episodic (red) and semantic (green) memory. Episodic memories integrate both sensory information (orange) and semantic content (green), while semantic memory extracts regularities across episodes. 
\textbf{b}: Schematic of bidirectional consolidation. Episodic-to-semantic consolidation transfers regularities from episodic memory to semantic memory, while semantic-to-episodic consolidation transfers the formed semantic representations to episodic memory.  
\textbf{c}: Circuit implementation of the model. Neocortex (CTX) stores semantic memories and Medial Temporal Lobe (MTL) stores episodic memories. MTL is split into neurons primarily receiving input from sensory cortex (MTL-sensory, coding for low-level sensory features) and neurons primarily receiving input from neocortex (MTL-semantic, coding for high-level abstract features). 
\textbf{d}: Mapping of cognitive consolidation processes onto replay mechanisms. Episodic replay corresponds to CTX$\leftarrow$MTL interactions (MTL-to-CTX connections trained during replay), while semantic replay corresponds to MTL-semantic$\leftarrow$CTX interactions (semantic assemblies in CTX driving concept representations in MTL-semantic).}
    \label{fig:diagram}
\end{figure}
\begin{comment}
\begin{tcolorbox}[colback=blue!5!white, colframe=blue!75!black, title=Technical Box - Input Generation]
\textbf{Intuition:} We generate compositional \emph{episodes} by combining one concept from each of two attributes. The different letters denote attributes (categories) and subscripts denote instances within each attribute. For example, $A_i$ could represent different visual cues and $B_j$ different auditory cues. An \emph{episode} is the pair $\{A_i,B_j\}$, and episodes are sampled from a joint distribution $p(A_i,B_j)$. Importantly, this joint distribution induces an asymmetric \emph{semantic structure} between concepts via conditional probabilities:
\[
\mathcal{SS}_{xy} \;=\; p(x \in e \mid y \in e),
\]
which we call ``semantic'' because it defines the meaning of a concept by the concepts that reliably co-occur with it. For instance, the semantics of \textit{summer} are given by concepts that are likely given \textit{summer} (sun, beach) and unlikely given \textit{summer} (snow, cold). 
\newline\newline
\textbf{Relevant details:}
\begin{itemize}
    \item \textbf{Sampling episodes within a day.} A day of sensory input has length $T_\textrm{day}$ timesteps and is partitioned into $T_\textrm{day}/T_\textrm{episode}$ blocks. At the start of each block, we sample one episode
    \[
    e=\{A_i,B_j\} \sim p(A_i,B_j),
    \]
    and hold its underlying concepts fixed for $T_\textrm{episode}$ timesteps.
    
    \item \textbf{Mapping an episode to a sensory activity pattern.} Each concept $c\in A\cup B$ is assigned a dedicated subpopulation of sensory neurons, $\mathrm{SEN}_c$ (Fig.~XXX). Given a sampled episode $e=\{A_i,B_j\}$, we generate an initial binary sensory vector $X^{\mathrm{SEN}}$ by activating exactly the neurons belonging to the two concept-specific subpopulations:
    \[
    X^{\mathrm{SEN}}_n =
    \begin{cases}
    1 & \text{if } n\in \mathrm{SEN}_{A_i}\ \text{or}\ n\in \mathrm{SEN}_{B_j},\\
    0 & \text{otherwise.}
    \end{cases}
    \]
    Operationally, $\mathrm{SEN}_{A_i}$ and $\mathrm{SEN}_{B_j}$ correspond to two concatenated ``blocks'' (one for attribute $A$, one for attribute $B$), with each block containing a fixed number of neurons per concept (Fig.~XXXX). This yields a structured, sparse sensory pattern whose active set explicitly encodes the current episode concepts.
    
    \item \textbf{Within-episode variability (swap noise).} To model variability across repeated presentations of the same episode, we inject controlled stochasticity at every timestep. Starting from the base pattern above, we uniformly sample $N_\textrm{swaps}$ active neurons and $N_\textrm{swaps}$ inactive neurons and flip their states (Fig.~XXXXX). This produces a noisy realization while preserving overall activity level (sparsity is unchanged because the number of $1\!\to\!0$ flips equals the number of $0\!\to\!1$ flips).
    
    \item \textbf{Temporal sequence.} Repeating the above over blocks produces a day-long tensor $X_\textrm{input}\in\{0,1\}^{(N_\textrm{days},T_\textrm{day}, N_\textrm{SEN}})$ in which (i) the episode identity changes every $T_\textrm{episode}$ timesteps and (ii) noise is independently resampled at each timestep.
\end{itemize}
\end{tcolorbox}
\begin{tcolorbox}[colback=blue!5!white, colframe=blue!75!black, title=Technical Box - Network Dynamics and Local Learning]
\textbf{Intuition:} All regions are composed of binary neurons operating under strong competition. At any given timestep, only a fixed number of neurons are allowed to be active. This \emph{K-winners-take-all} (K-WTA) mechanism enforces sparse, distributed representations and induces competition between neurons to represent incoming patterns. Learning is purely local and Hebbian, but combined with divisive normalization (homeostasis). The interaction between competition and normalization makes synaptic weights converge to scaled conditional probabilities, so that connectivity reflects statistical structure in the environment.
\newline
\textbf{Details: } Each region receives a pre-activation synaptic input:
\begin{equation}
    \hat{X}^\textrm{region} = W^{\textrm{region}\leftarrow\textrm{pre}}\cdot\hat{X}^\textrm{region}
\end{equation}
which is then mapped into binary activations via a top-K operation (therefore implementing \textit{K-winners-take-all} dynamics):
\begin{equation}
    X^\textrm{region} = \textrm{top-}K\big(\hat{X}^\textrm{region}\big)
\end{equation}
The basic form of learning is Hebbian: 
\begin{equation}
    \Delta W_{ij}^{\textrm{post}\leftarrow\textrm{pre}} = \lambda_{ij}X^\textrm{post}_iX^\textrm{pre}_j
\end{equation}
which is complemented with divisive normalization. When synaptic weights are divided by the sum of incoming weights $\sum_iW_{ij}$, we call this process \textit{incoming homeostasis}. Under this form of learning, the weights behave as an exponential moving average of conditional probabilities, converging to:
   \begin{equation}
   W_{ij}^{\textrm{pre}\leftarrow\textrm{post}, \; \textrm{incoming}}(\infty) \equiv \frac{W_\textrm{max-in}^{\textrm{pre}\leftarrow\textrm{post}}}{K^\textrm{pre}}p(X^\textrm{pre}_j = 1| X^\textrm{post}_i = 1)
 \end{equation}
and when these are divided by the sum of outgoing weights $\sum_iW_{ij}$, the fixed points are:
  \begin{equation}
     W_{ij}^{\textrm{pre}\leftarrow\textrm{post}, \; \textrm{outgoing}}(\infty) \equiv \frac{W_\textrm{max-out}^{\textrm{pre}\leftarrow\textrm{post}}}{K^\textrm{post}}p(X^\textrm{post}_i = 1| X^\textrm{pre}_j = 1)
 \end{equation}
 \textbf{Motivation:} The basic neuronal and learning dynamics are inspired by classic models of competitive learning and sparse coding CITE. The implementation details are heavily based on recent work 
\end{tcolorbox}
\end{comment}
\newtcolorbox[auto counter]{technicalbox}[2][]{
  colback=blue!5!white,
  colframe=blue!75!black,
  title=Technical Box~\thetcbcounter: #2,
  label=#1
}
\clearpage
\begin{technicalbox}[box:hbn]{Homeostatic Binary Networks}
All regions in the circuit model are implemented as Homeostatic Binary Networks, following \shortciteA{albesagonzalez_2025}. These networks are binary, with sparsity in each region being controlled by a \textit{K-winners-take-all} mechanism, implemented with a $top$-$K$ activation function. For further details and definitions of variables see Methods.
\newline\newline
\textbf{Neural dynamics.} Each region is described by a binary activity vector
\[
X^\mathrm{region}\in\{0,1\}^{N_\mathrm{region}},
\]
with exactly $K^\mathrm{region}$ active neurons at each timestep. Synaptic input is first integrated into a pre-activation variable
\[
\hat{X}^\mathrm{region}
=
W^{\mathrm{region}\leftarrow\mathrm{pre}}X^\mathrm{pre},
\]
which is then mapped into binary activity by competition:
\[
X^\mathrm{region}
=
\operatorname{top\text{-}K^\mathrm{region}}
\big(\hat{X}^\mathrm{region}\big).
\]
This mechanism causes enforces a competition across neurons to represent the current input.
\newline\newline
\textbf{Local learning.} Synapses are updated by a Hebbian rule,
\[
\Delta W_{ij}^{\mathrm{post}\leftarrow\mathrm{pre}}
=
\lambda_{ij}X_i^\mathrm{post}X_j^\mathrm{pre},
\]
followed by divisive normalization (homeostasis), which prevents runaway growth and keeps the total synaptic strength bounded.
\newline\newline
\textbf{Homeostatic regimes.} We use two forms of synaptic homeostasis. With \textit{incoming homeostasis}, weights are normalized so that the total input onto each postsynaptic neuron is bounded. With \textit{outgoing homeostasis}, weights are normalized so that the total output from each presynaptic neuron is bounded. Under these two regimes, synaptic weights converge to scaled conditional probabilities:
\[
W_{ij}^{\mathrm{pre}\leftarrow\mathrm{post},\,\mathrm{incoming}}(\infty)
\equiv
\frac{W_{\mathrm{max\text{-}in}}^{\mathrm{pre}\leftarrow\mathrm{post}}}{K^\mathrm{pre}}
\,p(X_j^\mathrm{pre}=1\mid X_i^\mathrm{post}=1),
\]
\[
W_{ij}^{\mathrm{pre}\leftarrow\mathrm{post},\,\mathrm{outgoing}}(\infty)
\equiv
\frac{W_{\mathrm{max\text{-}out}}^{\mathrm{pre}\leftarrow\mathrm{post}}}{K^\mathrm{post}}
\,p(X_i^\mathrm{post}=1\mid X_j^\mathrm{pre}=1).
\]
Thus, incoming homeostasis makes synapses approximate the average presynaptic pattern associated with a postsynaptic assembly, whereas outgoing homeostasis makes synapses reflect how strongly presynaptic neurons predict postsynaptic recruitment.
\newline\newline
\textbf{Functional role.} In recurrent connectivity, these dynamics implement pattern completion. In feedforward connectivity, they implement competitive dictionary learning, causing neurons to become selective to stable components of their input.
\end{technicalbox}
\begin{technicalbox}[box:replay]{Replay for Dictionary Learning}
Instead of viewing replay as perfect reinstatement of stored episodes, we propose it can be seen as a mechanism for dictionary learning, where a \textit{post} region learns relevant \textit{atoms} describing a \textit{pre} region based on presynaptic recurrent connections. The mechanism is as follows:
\newline\newline
\textbf{Generate a presynaptic pattern} from noise via attractor dynamics:
\[
X^\mathrm{pre}_{0} \sim \mathcal{N}(0,I_{N_\mathrm{pre}})
\quad \longmapsto \quad
X^\mathrm{pre}_{n+1}
=
\operatorname{top\text{-}K_R^{\mathrm{replay}}}
\left(
W^{\mathrm{pre}\leftarrow \mathrm{pre}} X^\mathrm{pre}_{n}
\right).
\]
\textbf{Project to postsynaptic layer}
\[
X^\mathrm{post}
=
\operatorname{top\text{-}K^\textrm{post}}
\left(
W^{\mathrm{post}\leftarrow\mathrm{pre}}X^\mathrm{pre}
+
b_\mathrm{IM}^\textrm{post}\mathrm{IM^\textrm{post}}
\right),
\]
\textbf{Update dictionary} For postsynaptic active neuron $i$ ($X^\mathrm{post}_i$ = 1).
\begin{itemize}
    \item If neuron $i$ is immature ($\textrm{IM}^\textrm{post}_i$ = 1): learn new atom
    \begin{itemize}
        \item Create one shot receptive field $W_{ij}^{\textrm{post}\leftarrow\textrm{pre}} \propto X^\textrm{pre}_j$ \\  
        $\longrightarrow$ neuron $i$ becomes maximally selective to $X_\textrm{pre}$
        \item Make neuron $i$ mature ($\textrm{IM}^\textrm{post}_i$ = 0)
    \end{itemize}
    \item If neuron $i$ is mature  ($\textrm{IM}^\textrm{post}_i$ = 0): update existing atom
    \begin{itemize}
        \item Receptive field is pushed in the direction of $X^\textrm{pre}$
        \[W_{ij}^{\textrm{post}\leftarrow\textrm{pre}} \longmapsto
p(X^\mathrm{pre}_j=1 \mid X^\mathrm{post}_i=1)\]
    \end{itemize}
\end{itemize}
\end{technicalbox}
\begin{comment}
\begin{tcolorbox}[colback=blue!5!white, colframe=blue!75!black, title=Technical Box - General Learning Algorithm]
\textbf{Intuition:} 
Learning alternates between two regimes: \emph{wake} (input-driven encoding) and \emph{sleep} (replay-driven abstraction). During wake, sensory input drives the network and updates recurrent synapses locally: MTL rapidly stores episodic traces, while CTX slowly accumulates statistical regularities. During sleep, internally generated replay transfers structure across regions. Depending on the training phase, replay progressively shifts from purely sensory episodic content to increasingly semantic and compositional structure.
\newline\newline
\textbf{Wake (online encoding):}
For each day, sensory input is processed sequentially:
\begin{enumerate}
    \item Sensory activity drives MTL-sensory and (initially random) MTL-semantic populations.
    \item MTL activity projects to CTX, where sparse competition selects a cortical representation.
    \item Recurrent synapses are updated locally:
    \begin{itemize}
        \item MTL recurrent connections store the current episode (rapid episodic encoding).
        \item CTX recurrent connections are updated with a slow learning rate, gradually capturing co-occurrence statistics (semantic structure).
    \end{itemize}
\end{enumerate}

At early stages, MTL-semantic activity is not yet driven by CTX and episodic traces contain only sensory content. After cortical structure has formed (Phase A), CTX begins to drive MTL-semantic, enriching episodic representations with semantic components. In later stages (Phase B), episodic traces contain both sensory and semantic structure, allowing higher-order regularities to be stored.
\newline\newline
\textbf{Sleep (replay-driven reorganization):}
Sleep consists of two replay loops:

\begin{itemize}
    \item \emph{Episodic replay (MTL $\rightarrow$ CTX):} 
    Recurrent MTL dynamics generate structured patterns which train CTX feed-forward synapses. Early in training, replay involves only sensory MTL content. After Phase B, replay draws from full MTL representations (sensory + semantic), and the semantic load can vary, enabling compositional abstraction.

    \item \emph{Semantic replay (CTX $\rightarrow$ MTL-semantic):} 
    Once cortical recurrent structure is established (after Phase A), CTX replay trains MTL-semantic representations, creating pointers to cortical statistical regularities.
\end{itemize}

Each replay step follows the same local rule: recurrent pattern completion in the \emph{pre} region, projection to the \emph{post} region, and Hebbian learning with homeostasis. The difference across phases lies only in which region replays to which.
\newline\newline
\textbf{Phase structure:}
\begin{itemize}
    \item \textbf{Phase A:} MTL stores sensory episodes during wake. Sleep replay trains CTX to extract sensory regularities.
    \item \textbf{Phase B:} CTX representations begin to drive MTL-semantic during wake. Sleep additionally transfers cortical structure back to MTL-semantic.
    \item \textbf{Late phase:} Episodic traces contain both sensory and semantic components. Replay operates over this richer space, enabling higher-order abstraction.
\end{itemize}
\newline
\textbf{Motivation:} There are some unusual choices in the general learning algorithm. For instance, the cortical generative model (embedded in cortico-cortical synapses) is not trained during sleep. Instead, sleep re-organizes cortical and medial temporal representations. A similar approach has been proposed in the wake-sleep library learning algorithm for program synthesis of DreamCoder CITE. Furthermore, \textit{episodic replay} can be mapped to the \textit{Abstraction} phase of DreamCoder, with both models creating new pointers that represent regularities present across episodes. In a similar fashion, \textit{semantic replay} can be compared to the corresponding \textit{Dreaming} (program recognition) sleep phase, with our network learning how to represent newly encoded episodes using the pointers discovered during episodic replay. Add summary regarding semantic replay biological plausibility. 
\end{tcolorbox}
\end{comment}
\newpage
\subsection*{A Neocortex - Medial Temporal Lobe model can form semantic memories via dictionary learning}
\subsubsection*{Neocortex forms a dictionary of reusable components via episodic replay}
We first consider episodic-to-semantic consolidation (Fig. 2a) as a form of dictionary learning carried out by neocortex. In dictionary learning, complex inputs are represented by combining a small number of reusable elements (typically called \textit{atoms}, Technical Box \ref{box:replay}). Here, we refer to these learned elements as \textit{cortical words}: reusable cortical representations that capture recurring parts of episodic experience. In this view, episodic replay in MTL does not simply reinstate full episodes. Rather, sparse replay proposes recurring subpatterns embedded across many episodes. CTX then either matches each replayed candidate to a previously learned cortical word, refining the corresponding receptive field, or treats it as sufficiently novel to recruit immature cortical neurons and establish a new receptive field. In this way, repeated episodic replay gradually builds a cortical dictionary of reusable semantic components (words).
\newline\newline
We first test whether CTX can implement such an algorithm in the model. To do so, we generate episodes by combining pairs of latent factors $(A_i,B_j)$, with $A_i \in \{A_1,\dots,A_5\}$ and $B_j \in \{B_1,\dots,B_5\}$ (Figs. \ref{fig:phase_a}b and \ref{fig:input_generation}, see Methods). Each latent factor (e.g. $A_2$) corresponds to a concept that gives rise to a semantically coherent set of sensory features, such as \textit{tomato} or \textit{cheese}. The distribution over pairings was structured: matched pairs, such as $(A_3,B_3)$, accounted for 50\% of episodes, whereas all mismatched pairings jointly accounted for the remaining 50\%. Thus, although every $A_i$ could in principle appear with every $B_j$, factors with matching indices co-occurred more often. For example, $A_1$ was more predictive of $B_1$ than $B_2$.
\newline\newline
During wake, episodic inputs are stored in MTL through Hebbian plasticity. During sleep, sparse pattern completion (Technical Box \ref{box:replay}) in MTL recovers sub-patterns shared across episodes (e.g. most active neurons correspond to $B_2$), rather than reinstating complete episodic states (Fig. \ref{fig:phase_a}c). These replayed subpatterns are projected to CTX through initially random CTX$\leftarrow$MTL connections (Fig. \ref{fig:phase_a}d, $t=0$).If the replayed pattern is not sufficiently close to existing cortical receptive fields (i.e. no CTX neuron is strongly tuned to the current MTL pattern), cortical activity is biased toward immature neurons (Methods). Because these neurons are more excitable, they are preferentially recruited, forming a one-shot cortical representation of the replayed MTL pattern (Fig. \ref{fig:phase_a}d, \textit{word proposal}, $t=200$). However, when the replayed pattern is sufficiently close to an existing cortical receptive field, the corresponding cortical word is reactivated and refined (Fig. \ref{fig:phase_a}d, \textit{word refinement}, from $t=200$ on). This process slowly tunes cortical neurons to stable, de-noised components of episodic experience (Fig. \ref{fig:phase_a}d,e).
\newline\newline
We further confirmed representational change in CTX by measuring the mutual information between cortical activity and the underlying \textit{simple concepts}, $A_i$ and $B_j$. Cortical activity rapidly became maximally informative about these latent factors, reaching a normalized mutual information of approximately 1 (Fig. \ref{fig:phase_a}f). \textit{Complex concepts} of the form $A_iB_j$ are considered in later sections. We also found that the probability of learning a cortical word increased monotonically with the probability that its corresponding simple concept appeared in an episode (Fig. \ref{fig:phase_a}g). Thus, this mechanism can be interpreted as a form of prioritizing replay for compositional consolidation: frequently recurring compositional elements are more likely to be consolidated into cortical words (semantic representations).
\subsubsection*{Cortical recurrent connectivity represents the semantic content of cortical words}
In our cognitive model (Fig. \ref{fig:diagram}a), we proposed that semantic memory contains contains information about the relatedness between concepts, which is defined based on their directional predictiveness (e.g. pizza is highly predictive of cheese). In the circuit model, when multiple cortical neurons become tuned to the same replayed pattern, they also become strongly recurrently connected, forming a cortical assembly for that word (Methods). This corresponds to the standard replay-mediated formation of recurrent cortical assemblies, in which re-activated neurons are bound into a stable attractor-like representation. However, during wake, slow plasticity further updates recurrent connections across the broader cortical population. Under Hebbian plasticity with outgoing homeostasis, these recurrent weights converge to the conditional firing structure between neurons \cite{albesagonzalez_2025}. Because cortical neurons are tuned to individual concepts, recurrent connectivity comes to reflect how predictive concepts are of one another (Figs. \ref{fig:phase_a}h and \ref{fig:circuit}b). Thus, while feedforward CTX$\leftarrow$MTL weights learn the dictionary entries---what each cortical word encodes--- recurrent CTX$\leftarrow$CTX connections learn the relations between entries---how each word is situated within the broader semantic structure. 
\newline\newline
Together, these results suggest that semantic memory can emerge from episodic replay through cortical dictionary learning. Replay extracts candidate words from overlapping episodes, CTX stabilizes them as reusable semantic units, and recurrent cortical connectivity organizes them according to their co-occurrence statistics. In the pizza example, cortical words encode ingredients such as cheese or tomato, whereas corticocortical connections encode the predictive relationships between ingredients.
\begin{figure}
    \centering
    \includegraphics[width=\linewidth]{Figures/Figure_2.png}
\caption{\textbf{Sparse episodic replay creates compositional representations in semantic memory.} \textbf{a}: Episodic-to-semantic consolidation via episodic replay. Cognitive model (left) and circuit implementation (middle/right). Sensory episodes (no semantics extracted yet) are stored in MTL-sensory, replayed sparsely during sleep to uncover overlapping sub-patterns, and then projected to CTX, where compositional semantic structure is extracted in CTX$\leftarrow$MTL (replay) and CTX$\leftarrow$CTX synapses (awake).  
\textbf{b}: Example neuronal activity during wake. Input episodes are sampled following a protocol as described in \shortciteA{albesagonzalez_2025}. To determine MTL-sensory activity, two different concepts (one $A_i \in A$ and one $B_j \in B$) are simultaneously sampled. 50\% of episodes are of form $(A_i, B_{j=i})$, and 50\% $(A_i, B_{j\neq i})$. This makes $A_i$ more semantically related to $B_i$ than the rest of $B$'s. 10 neurons in MTL-sensory code for each concept, and each episode lasts 5 timesteps. At each timestep, $N_\textrm{swaps}= 4$ random neurons that code for the present concepts are turned off, and  $N_\textrm{swaps}= 4$ random neurons that code for non-present concepts are turned on (this introduces noise with a consistent signal-to-noise ratio). CTX activity is driven by random initial CTX$\leftarrow$MTL-sensory connections.
\textbf{c}: Example neuronal activity (one timestep) during wake (left) and episodic replay (right). Sparse replay in MTL-sensory reveals concept-specific sub-patterns (for example, here $B_2$).}
    \label{fig:phase_a}
\end{figure}
\begin{figure}[t]
  \ContinuedFloat
  \captionsetup{list=off} % don't add a new list-of-figures entry
  \caption{\textit{(continued)} \textbf{d}: Evolution over time of the receptive field (CTX$\leftarrow$MTL synapses) of an example neuron. Right: weight dynamics for all synapses of the same neuron. Left: Heatmaps show weights associated to each presynaptic MTL-sensory neuron rearranged in a rectangular pattern according to the concept they code for (each row represents weights from neurons coding for a particular concept in MTL-sensory). Initially random connectivity ($t = 0$) is one-shot driven to represent a replayed presynaptic pattern ($t=200$), and neuron 60 goes from immature to mature. Then, slow Hebbian and homeostatic plasticity drive statistical learning, uncovering receptive fields that represent semantic components ($t = 200-90000$) \textbf{e}: Comparison of early (before learning) and late (after learning) cortical activity, given the same sensory input. In the late pahse, CTX neurons now respond selectively to semantic components $A_i$ and $B_j$. \textbf{e}: Evolution with time of mutual information of cortical activity with simple concepts (e.g. $A_1$). Average over concepts. \textbf{g}: For different marginal probabilities of a concept to be present in an episode (horizontal axis), probability of a receptive field selective to that concept to be added in CTX during one session of sleep (vertical axis). \textbf{h}: For CTX$\leftarrow$CTX connections, relationship between synaptic strength and semantic structure (right). Recurrent connectivity captures conditional firing probabilities across concepts.}
\end{figure}
\newpage
\subsection*{Semantic replay can transfer the learned vocabulary into episodic memory}
We next asked whether spontaneous cortical activity can transfer the learned cortical dictionary back into episodic memory (semantic-to-episodic consolidation; Fig. \ref{fig:phase_b}a, left). We refer to this mechanism as \textit{semantic replay} because it parallels episodic replay: it is not directly driven by sensory input, and activity patterns are recovered through attractor dynamics from random initial states. However, here the replayed content comes from recurrent connections representing \textit{semantic} information, as opposed to episodic. In the pizza example, once ingredients such as tomato and cheese have been consolidated as cortical words, semantic replay can imprint these words onto MTL-semantic. Thus, before semantic-to-episodic consolidation, an episode containing tomato sauce may be represented only through sensory features. After consolidation, MTL-semantic additionally contains neurons that explicitly signal the presence of tomato, analogous to MTL concept cells. At the circuit level, this corresponds to MTL-semantic neurons developing receptive fields that are selective for cortical assemblies representing individual concepts (Fig. \ref{fig:phase_b}a, right).
\newline\newline
To support this process, we now equip the simulated model with semantic replay (Algorithm \ref{alg:sleep}). Semantic replay initializes CTX from random activity and allows recurrent cortical dynamics to recover learned semantic attractors. As in the previous section, sparsity is increased to isolate patterns corresponding to individual concepts (Fig. \ref{fig:phase_b}b). These patterns are then projected to MTL-semantic through initially random connections. The rationale mirrors that of episodic replay: increasing sparsity produces non-overlapping patterns that can be efficiently learned via competitive learning. As a result, receptive fields in MTL-semantic come to represent individual concepts identified by CTX (Fig. \ref{fig:phase_b}d). Consequently, while MTL-sensory exhibits a mixed selectivity to different concepts, MTL-semantic neurons show an absolute selectivity to these concepts (Fig. \ref{fig:phase_b}d, e), with mutual information becoming maximal rapidly after semantic replay is initiated (Fig. \ref{fig:phase_b}f).
\newline\newline
Once MTL-semantic has acquired concept-selective receptive fields, replay can engage both MTL-sensory and MTL-semantic populations through recurrent MTL connectivity (Fig. \ref{fig:phase_b_supp}). As a result, subsequent replay events can contain both sensory and semantic information, inherited from CTX. While semantic representations are identical in content in CTX and MTL, they serve different purposes. In CTX, recurrent connectivity between semantic assemblies stores long-term perdictiveness across concepts. However, in MTL-semantic, recurrent connectivity stores the meaning associated to the different episodes stored during a day. Extending replay in MTL to joint sensory-semantic episodes expands the previous CTX$\leftarrow$MTL-sensory connections to form cortical receptive fields that jointly respond to MTL-sensory and MTL-semantic patterns (Fig. \ref{fig:circuit}d). These connections will be shown to be crucial for consolidating complex combinations in the following sections.
\newline\newline
Overall, these results show that semantic replay implements semantic-to-episodic consolidation in the circuit model, as proposed in the cognitive model. This provides a mechanistic account of how learned semantic representations can be transferred from neocortex to MTL, giving rise to concept-like cells \shortcite{quiroga_2005,quiroga_2012}, and allowing CTX and MTL to semantically index each other via a shared dictionary.
\begin{figure}
    \centering
    \includegraphics[width=\linewidth]{Figures/Figure_3.png}
\caption{\textbf{Semantic replay forms abstract representations in episodic memory.} \textbf{a}: Semantic-to-episodic consolidation via semantic replay. Cognitive model (left) and circuit implementation (right). During semantic replay, a random CTX pattern converges to a concept-specific attractor, which is then projected to MTL-semantic, creating semantic-specific receptive fields in MTL-semantic and enabling episodic traces to incorporate semantic information. \textbf{b}: Example neuronal activity before semantic replay starts. CTX has organized its activity around semantic components, but MTL-semantic is still a random projection of cortical activity.
\textbf{c}: Example neuronal activity during semantic replay. A random initial CTX state evolves into a complete pattern coding for a single concept, which drives activity in MTL-semantic. A combination of one-shot and slow learning again make neurons that are randomly slightly more sensitive to cortical patterns selective to cortical semantic representations.
\textbf{d}: Same as Fig. \ref{fig:phase_b}b, but after learning. Example neuronal activity before semantic replay starts. MTL-semantic activity is now also organized around semantic components.
\textbf{e}: Selectivity of neurons in MTL-sensory and MTL-semantic population  (*** $p < 0.001$).
\textbf{f}: Evolution of mutual information between different MTL-subregions and simple concepts (averaged across concepts).
}
\label{fig:phase_b}
\end{figure}
\newpage
\subsection*{Variable sparsity during replay controls the complexity of the vocabulary learned}
So far, episodic replay has been constrained to high sparsity, a regime that decorrelates episodic subpatterns and promotes the extraction of \textit{simple concepts} (e.g. tomato or cheese). However, if replay is allowed to recruit a larger number of active neurons (decreasing sparsity), the network can in principle recover patterns corresponding to higher-order combinations of concepts. For example, replay may recover conjunctions of concepts such as \textit{crust}, \textit{cheese}, and \textit{tomato} by allowing a less sparse pattern to emerge (Fig. \ref{fig:higher-order}a).
\newline\newline
To test this idea, we extended the replay dynamics to allow variable sparsity. We define the \textit{semantic load} as the ratio between the number of active neurons during replay and the number active during wakeful episodic encoding. Intuitively, semantic load controls how many concepts are simultaneously expressed in a replayed pattern. In the current setting, replay events have a semantic load of either 1 or 2 (Fig. \ref{fig:higher-order}b). Consequently, MTL attractor dynamics recover a pattern associated  with either single concept (e.g. $A_1$) or a conjunction of two concepts (such as $A_1B_2$), depending on how many neurons are allowed to remain active. The resulting MTL activity is then projected to CTX, where a separate cortical subregion learns these higher-order combinations.
\newline\newline
Under this regime, replay with a semantic load 2 induces a competition between neurons coding for specific combinations $(A_i,B_j)$, that is, complete episodic conjunctions rather than isolated concepts. Over time, CTX develops selective receptive fields for these combinations (Figs. \ref{fig:higher-order}b-d, \textit{complex concepts}, and Fig. \ref{fig:circuit}d). As a result, cortical activity during wake contains neurons tuned to each $A_i$, to each $B_j$, and to specific conjunctions $(A_i,B_j)$ (Fig. \ref{fig:higher-order}c). In this way, cortical dictionary learning extends from simple cortical words (\textit{atoms} represent simple concepts) to higher-order compositional words (\textit{atoms} represent complex concepts).
\newline\newline
When examining CTX neurons selective to complex concepts, we find that learning reorganizes activity into a conjunctive code, with the region representing at each timestep representing an abstracted individual episode $A_iB_j$ (Fig. \ref{fig:higher-order}d). Consistently, the mutual information between these neurons and complex latent factors of the form $A_iB_j$ increased sharply once variable sparsity was introduced (Fig. \ref{fig:higher-order}e). By contrast, the normalized mutual information of CTX simple-concept neurons and MTL-sensory activity with these complex latent factors remained notably below maximal (Fig. \ref{fig:higher-order}e). Similarly to the case of simple concept voabulary learning, the formation of new receptive fields was favoured by an increased likelihood of the corresponding episode to be sampled (Fig. \ref{fig:higher-order}f).
\newline\newline
This transformation is also reflected in corticocortical connectivity. As higher-order cortical words form, CTX recurrent synapses encode the relational structure between simple concepts and their conjunctions (Fig. \ref{fig:circuit}e). It should be noted the asymmetry of this structure is inherited by the connectivity matrix: a conjunction such as $A_1B_1$ strongly predicts its constituent $A_1$, whereas $A_1$ only weakly predicts the specific conjunction $A_1B_1$ (Fig. \ref{fig:circuit}e). Consequently, recurrent connectivity can encode directional predictive relationships, overcoming a limitation of many Hebbian models that are restricted to symmetric relationships such as correlation.
\newline\newline
Together, these results show that replay sparsity controls the level of abstraction learned by the cortical dictionary. Highly sparse replay favours isolated concepts, whereas less sparse replay allows CTX to consolidate higher-order cortical words corresponding to concept conjunctions. In the pizza example, replay may first extract individual ingredients such as \textit{cheese}, \textit{tomato}, or \textit{crust}, and later consolidate pizza as their conjunction, capturing increasingly episodic-like semantic structure.
\begin{figure}[h!]
    \centering
    \includegraphics[width=0.95\linewidth]{Figures/Figure_4.png}
\caption{\textbf{Learning higher-order concepts based on simpler representations is facilitated by the presence semantic representations during replay.}}
    \label{fig:higher-order}
\end{figure}
\begin{figure}[h!]
  \ContinuedFloat
  \captionsetup{list=off} % don't add a new list-of-figures entry
  \caption{\textit{(continued)} \textbf{a}: Cognitive model (left) and circuit implementation (right) of higher-order concept learning. Multiple sensory–semantic episodes are stored in MTL. Variable sparsity during replay (semantic load) allows the recovery of higher-order associations, which are then consolidated into CTX$\leftarrow$MTL and  CTX$\leftarrow$CTX connections.\textbf{b}: Example activity during replay with semantic load 1 (top) or 2 (bottom). When one concept is replayed, CTX develops selective coding for single concepts. When two concepts are replayed simultaneously, CTX develops receptive fields for their combinations, supporting higher-order abstractions.  
\textbf{c}: Example activity during wake after training. CTX contains neurons selective to individual concepts $A_i$ and $B_j$, as well as to combinations ($A_i, B_j$).  
\textbf{d}: Learning reorganization of CTX activity (complex concepts selective neurons). Conjunctive representations form duplicated representations of $B_j$ according depending on the $A_i$ context (can also be interpreted in reverse).
\textbf{e}: Evolution of mutual information with complex (e.g. $A_1B_3$) across different regions. Average across  complex concepts. \textbf{f}: Same as Fig. \ref{fig:phase_a}g, but for complex concepts.}
\end{figure}
\subsection*{The model reproduces key neural and synaptic experimental signatures}
We first tested whether semantic structure in the model reshapes medial temporal coding in a manner consistent with the defining properties of concept cells, namely concept-level representations that generalize across contexts (Fig. \ref{fig:validation_circuit}a). A central commitment of the model is that semantic abstraction is represented within the episodic circuit itself, rather than emerging only at a downstream cortical readout. This predicts that MTL populations should support concept-level generalization across contexts. Consistent with this, concept information was more readily decoded across held-out contexts from MTL-semantic than from MTL-sensory (Fig. \ref{fig:validation_circuit}b1). At the single-neuron level, MTL-semantic units also showed stronger latent selectivity (Fig. \ref{fig:validation_circuit}b2) and carried more information about their preferred concept (Fig. \ref{fig:validation_circuit}). Together, these results are consistent with the idea that semantic structure emerges within episodic circuitry and gives rise to concept-like neural codes.
\newline\newline
We next asked whether consolidation reshapes representational geometry in a manner consistent with the model's overlap-based replay mechanism. Because replay is structured by shared latent factors, the model predicts that memories sharing latent structure should become more similar over consolidation, both in hippocampus and in cortex.  A related effect was reported by \citeA{tompary_2017} in an object-scene paradigm, in which multiple object memories shared a smaller number of scene associations. The authors measured representational similarity during object-cued recall before and after sleep for objects associated with overlapping or non-overlapping scenes. Importantly, this task maps naturally onto our episode generation protocol, with $A_i$ corresponding to objects and $B_j$ to scenes (Fig. \ref{fig:validation_circuit}c). Consistent with the experimental findings, consolidation increased the overlap index in cortex, such that objects sharing a scene became more similarly represented after sleep than before sleep (Fig. 5d1). A corresponding increase was also observed in MTL (Fig. 5d2).
\newline\newline
Finally, we asked whether the particular neural and synaptic changes generated by the model could be related to experimentally described circuit motifs. Recent work \shortcite{lee_2023} characterized synaptic and representational changes during consolidation in a contextual fear-conditioning task, offering a particularly useful benchmark for our model because it further constrains the underlying circuit mechanisms. These experiments focused on interactions between the medial prefrontal cortex (mPFC) and basolateral amygdala (BLA), which we map onto the neocortex (CTX) and medial temporal lobe (MTL), respectively, in our model. To capture the experimental paradigm within our framework, we interpret the latent pair $(A_i,B_j)$ in a contextual fear-conditioning framework: different $A_i$ represent sensory descriptions of context, with one context serving as the conditioned stimulus (CS), whereas $B_j$ represents internal valence states, including fear as the unconditioned stimulus (US; $B_1$) (Fig. \ref{fig:validation_circuit}f1). Conditioning is then modeled by presenting the CS alone, followed within the same day by its pairing with the fear state.
\newline\newline
\citeA{lee_2023} examined how remote contextual fear memories are represented by distinguishing neurons active during conditioning from those active during later recall. In their terminology, \emph{engram neurons} are neurons recruited during the original fear-conditioning episode, whereas \emph{recall neurons} are neurons recruited during remote memory retrieval. These populations only partially overlap: some recall neurons are reactivated engram neurons, whereas others are newly recruited at recall. The model reproduces this reorganization of cortical recruitment (Fig. \ref{fig:validation_circuit}f2). In our framework, this follows naturally from cortical dictionary learning: the initial episode is first encoded using the pre-existing cortical vocabulary, recruits a dedicated cortical representation of the context-fear conjunction. As a result, remote recall can engage neurons that were not part of the original encoding ensemble, reproducing the partial overlap between engram and recall populations.
\newline\newline
When examining the connectivity motifs, the model developed stronger recurrent synapses among cortical engram neurons than between engram and non-engram neurons, reproducing the selective strengthening of engram-defined cortical connectivity after consolidation (Fig. \ref{fig:validation_circuit}g1). The model likewise reproduced stronger projections from cortical engram neurons to medial temporal engram neurons than from non-engram cortical neurons (Fig. \ref{fig:validation_circuit}g2), as well as stronger inputs from recall-recruited cortical neurons onto cortical engram cells than from non-recall neurons (Fig. \ref{fig:validation_circuit}g3). In the model, the strengthening of CTX$\rightarrow$MTL projections arises from semantic replay, which repeatedly reactivates cortical semantic assemblies and transfers their structure to MTL. Thus, the experimentally observed increase from the newly formed cortical engram to BLA is consistent with the prediction that semantic knowledge is progressively incorporated into episodic representations through cortex-to-MTL teaching.
\newline\newline
Lee et al. \cite{lee_2023} next asked whether extinction erases the cortical synaptic engram or instead weakens its expression. In particular, they compared engram-to-engram and engram-to-non-engram cortical synapses after fear conditioning followed by extinction, and found that the strong synaptic contrast observed after consolidation was no longer clearly maintained. The model reproduces this qualitative reduction in synaptic separation after extinction (Fig. \ref{fig:validation_circuit}i-j). In our framework, extinction does not erase the cortical representation formed during consolidation. Instead, it weakens the effective coupling from the context representation to the context-fear assembly, reducing the probability that the fear representation is reactivated during recall. As a result, extinction suppresses recall expression without fully erasing the underlying representational scaffold, providing a natural explanation for the enhanced re-learning (savings) observed when conditioning is reintroduced after an initial conditioning-extinction sequence (Fig. \ref{fig:validation_circuit}k).
\newline\newline
Together, these results link the core mechanisms of the model to experimentally observed neural and synaptic signatures. The cortical dictionary learning is consistent with consolidation-related increases in overlap of conceptually related memories \shortcite{tompary_2017} and the reorganization of cortical recuruitment observed between encoding and remote recall \shortcite{lee_2023}. Similarly, while standard consolidation theories predict catastrophic forgetting following extinction, our model explains savings in re-consolidation as an update of the semantic structure across previously existing cortical entries (embedded in recurrent cortical connections). In turn, the transfer of semantic representations into the episodic circuit is supported by the formation of concept cells \shortcite{quiroga_2005, quiroga_2012} and the fact that consolidation promotes overlap across conceptually similar memories not only in cortex, but also in hippocampus \shortcite{tompary_2017}. Finally, the strengthening corticofugal (MTL$\leftarrow$CTX) connections from cortical recurrent assemblies recruited during consolidation \shortcite{lee_2023} is consistent with the idea that spontaneous cortical activity can shape representations in downstream subcortical areas (\textit{semantic replay}).
\begin{figure}
    \centering
    \includegraphics[width=0.9\linewidth]{Figures/Figure_5.png}
    \caption{\textbf{The model reproduces key neural and synaptic signatures associated with semantic memory consolidation and concept representations.}}
    \label{fig:validation_circuit}
\end{figure}
\begin{figure}[h!]
  \ContinuedFloat
  \captionsetup{list=off} % don't add a new list-of-figures entry
  \caption{\textit{(continued)} 
\textbf{a}: Example of a human medial temporal lobe concept cell selective for the concept “Diego Maradona”, responding across different visual instances while remaining largely inactive to other familiar individuals. Adapted from Quiroga et al. (2005).
\textbf{b1}: Cross-context decoding of latent concepts from MTL activity. Concept identity can be decoded across held-out sensory contexts more accurately from MTL-semantic than from MTL-sensory representations, indicating stronger context-invariant coding.
\textbf{b2}: Maximum latent selectivity of individual neurons. MTL-semantic neurons exhibit higher selectivity for underlying latent concepts than MTL-sensory neurons.
\textbf{b3}: Mutual information between neuronal activity and concept identity. MTL-semantic neurons carry substantially more information about their preferred concept, consistent with concept-like representations.
\textbf{c}: Experimental paradigm adapted from Tompary and Davachi (2017). Multiple object memories share overlapping scene associations, allowing the assessment of consolidation-induced changes in representational similarity between related memories.
\textbf{d1}: Representational overlap in neocortex (CTX) before and after consolidation. Memories sharing latent structure become more similar after consolidation, reproducing the increased overlap observed experimentally for related memories.
\textbf{d2}: Representational overlap in the medial temporal lobe (MTL) before and after consolidation. Consolidation similarly increases overlap between memories sharing common latent factors within episodic circuitry.
\textbf{e}: Contextual fear-conditioning paradigm adapted from Lee et al. (2023). A conditioned context (CS) is paired with a fear-related unconditioned stimulus (US), enabling analysis of consolidation-induced circuit reorganization.
\textbf{f1}: Example activity used in the fear-conditioning simulation. Context representations are paired with a fear latent factor during conditioning.
\textbf{f2}: Cortical population activity during encoding (engram) and remote recall. Consolidation reorganizes cortical recruitment, producing only partial overlap between neurons active during encoding and those active during subsequent recall.
\textbf{g1}: Recurrent cortical connectivity after consolidation. Synapses between cortical engram neurons (E–E) become stronger than connections involving non-engram neurons, reproducing experimentally observed engram-specific strengthening.
\textbf{g2}: Corticofugal connectivity from cortical engram neurons to MTL engram neurons. Consolidation strengthens projections from cortical assemblies recruited during learning onto episodic-memory populations.
\textbf{g3}: Connectivity from recall-recruited cortical neurons onto cortical engram neurons. Newly recruited recall neurons exhibit stronger connections to the original cortical engram than non-recall neurons.
\textbf{h}: Extinction-learning paradigm adapted from Lee et al. (2023). Fear conditioning is followed by repeated extinction sessions in which the conditioned context is presented without reinforcement.
\textbf{i}: Recurrent cortical connectivity after extinction. The contrast between engram-to-engram and engram-to-non-engram synaptic strengths is reduced following extinction, consistent with experimental observations.
\textbf{j}: Extinction suppresses expression of the fear representation across extinction sessions.
\textbf{k}: Reacquisition after extinction. Previously extinguished memories are relearned faster than naive memories, demonstrating savings arising from the persistence of the underlying semantic and cortical representations.}
\end{figure}
\newpage
\subsection*{Normative benefits of semantic representations in episodic memory}
The emergence of semantic representations in the model's Medial Temporal Lobe allows us to ask what role they play in episodic memory and consolidation. The aim of this section is to clarify what are the normative benefits of embedding semantic structure within episodic memory. This question is particularly important in light of a longstanding tension between theory and experiment. On the one hand, classical attractor theories of episodic memory predict optimal performance for orthogonalized, uncorrelated patterns. On the other hand, experimental data suggest that episodic representations are far from unstructured, and instead contain clear semantic organization shared across memories. Furthermore, rather than being detrimental, semantic information can facilitate recall and is usually more easily remembered than sensory detail. Our model provides a framework to examine why such structure might be beneficial, and what computational trade-offs it introduces for recall, replay, and consolidation.
\newline\newline
To isolate the contribution of semantic representations, we compared three network modes (Fig. \ref{fig:normative}a). In the \textit{semantics present} condition (Fig. \ref{fig:normative}a, blue), the model operates normally and MTL-semantic contains concept-level representations learned from CTX. In the \textit{semantics random} condition (Fig. \ref{fig:normative}a, gray), MTL-semantic remains active but its activity is replaced by random patterns at each timestep, controlling for the presence of additional neural activity without semantic content. Finally, in the \textit{semantics absent} condition (Fig. \ref{fig:normative}a, red), MTL-semantic is excluded from MTL attractor dynamics altogether, corresponding to a purely sensory episodic system. Together, these conditions allow us to distinguish the effects of semantic information itself from the effects of simply increasing the size or activity of the episodic circuit.
\subsubsection*{Semantic representations in episodic memory increase memory capacity when episodes are semantically in-distribution}
We often recall events that we can semantically comprehend much better than those we cannot \shortcite{murray_1976, lin_2021}. This typically leads to an effective amnesia of that which we cannot semantically process (as introduced in \shortcite{tse_2007}). In order to answer what mechanism might lead to such effect, we first confirm that it in fact present in our model.
\newline\newline
We therefore begin by establishing a baseline in which no semantic structure is present. Specifically, we compare recall of purely random sparse patterns in our medial temporal lobe network against a classical sparse Hopfield network with matched size and sparsity (Fig. \ref{fig:normative}b). Under these conditions, our network showed only a minor advantage at low memory loads, followed by a steeper recall cliff as load increased. Thus, the model exhibits the standard overload regime expected for associative storage of random, unstructured patterns.
\newline\newline
We then examined recall as the number of stored patterns increased when episodes were sampled from our standard structured distribution (Fig. \ref{fig:normative}c), across different levels of input Signal-to-Noise Ratio (SNR). One striking difference with unstructured pattern storage, common across all networks, is the presence of a U-shaped learning curve. Performance initially decreased as more episodes were stored, but later recovered as repeated exposure to the same underlying structure allowed recurrent weights to average out sensory noise.
\newline\newline
Notwithstanding that, this recovery was weaker when semantic representations were absent or randomly assigned. Without semantic representations, the dip was more pronounced, whereas networks with aligned semantic representations retained higher recall. The advantage of semantic representations (measured as the difference across networks) also increased with sensory noise (Fig. \ref{fig:normative_supp}a). Mechanistically, MTL-semantic provides an additional, structured population that compresses the episode into concept-level variables. Although this does not increase the information content of the full episode, it increases the number of neurons available to represent it and provides an error-correcting scaffold for MTL-sensory activity. As a result, semantic representations improve not only semantic recall but also recall of the sensory component.
\newline\newline
Finally, within the \textit{semantics present} network, recall from MTL-semantic was substantially higher than recall from MTL-sensory. This is expected because MTL-semantic represents the lower-dimensional conceptual structure of the episode using its own dedicated neural population, whereas MTL-sensory must preserve more variable perceptual detail. This provides a circuit-level account of why the semantic content of memories is often more robustly recalled than their sensory details, which are more prone to being distorted or forgotten.
\newline\newline
Together, these results suggest that semantic structure optimizes the recall capacity of episodic memory by finding a compressed code that effectively amplifies the signal of the episodic conceptual structure. When stored patterns are random, the model behaves like a standard associative memory and eventually fails through overload. When episodes are semantically in-distribution, however, overlap is no longer purely interfering: it reflects reusable structure that MTL-semantic can encode compactly and use to stabilize recall. Thus, semantic representations do not generically increase memory capacity; rather, they increase effective capacity when new episodes can be interpreted through a learned semantic vocabulary.
\subsubsection*{Semantic representations in episodic memory induce a generalization bias in replayed and consolidated representations}
Following up on new theoretical studies proposing that replay optimizes generalization in Complementary Learning Sytems (GO-CLS, \citeA{sun_2023}), we hypothesize semantic representations can play a role in the generalization of rerpresentations consolidated via replay.
\newline\newline
First, we investigated whether the presence of semantic content during episodic encoding has an effect on the generalization of newly learned concepts. To do so, we pre-trained network that learned only simple concepts (by limiting semantic load during replay), and created 3 copies: \textit{semantics present}, \textit{semantics random}, and \textit{semantics absent} (Fig. \ref{fig:normative}a). Then, we increased the maximum number of neurons active during replay (semantic load), allowing it to consolidate complex concepts that were combinations on the previously learned simple concepts (e.g. $A_1B_2$ is a combination of $A_1$ and $B_2$). We then measured the accuracy of CTX complex representations at test time (compositional generalization).  We observed that as Signal-to-Noise Ratio (SNR) decreases, there is a higher gap between the generalization of cortical representations between networks with present and absent (or random) semantic representations (Fig.\ref{fig:normative}b). This gap persists for intermediate level of noise, until the accuracy of both converge to chance.
\newline\newline
To understand the origin of this cortical generalization gap, we next examined the replayed representations themselves. The first measure we obtained was the Signal-to-Noise Ratio of the replayed pattern. This was done by inferring what SNR in the input generation would generate a pattern as the one replayed (only looking at MTL-sensory). Semantic representations consistently increase the SNR associated to the replayed patterns, coinciding with the region in which there is a generalization gap (Figs. \ref{fig:normative}b and \ref{fig:normative}c). However, when explicitly testing the generalization of replayed MTL-sensory against awake patterns, the effect seems to not explain the gap in full (Figs. \ref{fig:normative}d and \ref{fig:normative_supp}b). That is, although the initially formed receptive field was slightly better at identifying future semantically equivalent episodes, this did not fully correlate with the generalization gap during testing. Then, we instead measured the Replayed to Replayed generalization: (how much more similar are replayed patterns that represent the same episode vs those that do not?, Fig. \ref{fig:normative}e). This turns out to explain much better the generalization gap (Fig. \ref{fig:normative}f, right). 
\newline\newline
These results can be interpreted in light of two slightly different mechanisms potentially driving better generalization after learning: when an episodic pattern is replayed, MTL is projected to CTX, and the activated cortical neurons update their MTL receptive fields in the direction of the replayed medial temporal pattern. Therefore, the formation of cortical representations during replay depends on two factors: (i) the similarity of the replayed MTL-sensory pattern with semantically equivalent future inputs and (ii) that the \textit{right} cortical neurons are activated, which depends on the generalization within replayed patterns. The second being more important for generalization after training is consistent with the idea that cortical representations are not determined by any single replay event, but instead emerge gradually through the accumulation of many replay-driven updates under a small learning rate. As a consequence, even if semantic structure also makes individual replayed MTL-sensory patterns less noisy, the dominant factor for cortical generalization is whether the \textit{appropriate} cortical neurons are consistently recruited across replay events. In other words, what matters most is not only that replay is less noisy, but that replay repeatedly drives learning onto the correct cortical assemblies.
\newline\newline
These results suggest that semantic representations promote cortical generalization not only by denoising replayed medial temporal patterns, but by providing a low-dimensional scaffold that remains tightly coordinated with cortex and helps constrain which cortical neurons are activated during replay. Together, our findings support the view that semantics shape replay into a teaching signal that is inherently more suitable for forming abstract, robust, and generalizable cortical representations.
\begin{figure}
    \centering
    \includegraphics[width=0.95\linewidth]{Figures/Figure_6.png}
    \caption{\textbf{Semantic representations in episodic memory improve recall capacity and promote generalization during replay and consolidation.}
\textbf{a}: Three network conditions used throughout the analysis. In the \emph{semantics present} condition, MTL-semantic contains concept representations learned from CTX. In the \emph{semantics random} condition, MTL-semantic activity is replaced by random patterns with matched activity levels. In the \emph{semantics absent} condition, MTL-semantic is removed from episodic dynamics, yielding a purely sensory episodic memory system.
\textbf{b}: Compositional generalization of newly learned cortical representations as a function of sensory input signal-to-noise ratio (SNR). Networks containing semantic representations in episodic memory show improved generalization, particularly at intermediate noise levels.
\textbf{c}: Effective signal-to-noise ratio of replayed MTL-sensory patterns as a function of input SNR. Semantic representations increase the quality of replayed patterns, indicating a denoising effect during attractor dynamics.
\textbf{d}: Replay-to-awake generalization. Similarity between replayed MTL representations and corresponding awake representations as a function of input SNR. Semantic representations moderately improve the fidelity of replayed episodic patterns.
\textbf{e}: Replay-to-replay generalization. Similarity between replay events representing the same latent episode relative to different episodes. Semantic representations substantially increase consistency across replay events, providing a more reliable learning signal for cortical consolidation.
\textbf{f}: Relationship between cortical test-time generalization and replay statistics. Left: compositional generalization correlates with replay-to-awake generalization. Right: compositional generalization is more strongly predicted by replay-to-replay generalization, suggesting that reliable recruitment of the same cortical assemblies across replay events is a key determinant of abstraction and generalization.
\textbf{g}: Episodic recall performance as a function of memory load under different levels of sensory noise (SNR = 2.25, 0.25, and 0.09). Dashed lines indicate recall from MTL-semantic representations, solid lines indicate recall from MTL-sensory representations, and the black curve indicates the cumulative number of distinct concepts encountered. Semantic representations increase effective memory capacity and improve recall of semantically structured episodes, particularly under noisy conditions.}
    \label{fig:normative}
\end{figure}
\subsection*{Semantic representations in the episodic memory circuit explain a range of behavioral phenomena}
We finally asked whether the properties of episodic attractor dynamics characterized above could explain a broader range of behavioral phenomena (Fig. \ref{fig:behavioural}). 
\newline\newline
In Fig. \ref{fig:normative}d, we showed how semantic MTL representations structure episodic memory to benefit recall and consolidation dynamics. This should be consistent with behavioural studies showing that semantically interpretable episodes are more faithfully recalled. For example, when comparing actual words with scrambled letters \shortcite{murray_1976}, or intact images with phase-scrambled ones \shortcite{lin_2021} (Fig. \ref{fig:behavioural}a). This result stands at the core of the abovementioned tension between theory (patterns should be sepparated as much as possible), and experiments (structured episodic overlap is benefitial for recall). To reproduce this effect, we compared the frozen network's recall of standard patterns (\textit{semantics intact} to one storing fixed permutations of the same patterns \textit{semantics scrambled}) (see Methods). The recall performance of patterns that could be semantically interpreted was significantly higher than that of semantically scrambled versions (Fig. \ref{fig:behavioural}c), even though the generated episodes were equivalent except for sensory neuronal identity.
\newline\newline
Another key difference of the model with previous accounts of semantic consolidation is that the episodic memory system performs an initial generalization step, finding structure across the different stored episodes in the medial temporal lobe. This was confirmed in Fig. \ref{fig:normative}e, where replayed activity was consistently de-noised with respect to sensory input. This predicts that enforcing a consistent overlapping component across all episodes benefits the formation of that particular semantic representation, which has been previously shown in humans in the literature of \textit{blocked} versus \textit{interleaved} learning.  Importantly, this result goes against standard neural network accounts of learning, and a circuit account using only local plasticity is yet to be proposed. We therefore tested a similar protocol in our network, training from scratch two networks: one with an initial blocked (all episodes in a day contain a fixed $A_i$ or $B_j$) phase, and one with an initial \textit{interleaved} phase, where every episode was equally likely (see Methods). Across noise levels, blocked learning showed a significant advantage over interleaved learning (Fig. \ref{fig:behavioural}d), replicating previous results. Furthermore, this was due to replayed patterns being more aligned with neural activity prototypes corresponding to the component fixed during each block (Fig. \ref{fig:behavioural}X). The model suggests that the mechanism via which humans benefit from learning in a block manner is that it promotes the consolidation of cortical representations representing each independent block, which remains hard to explain using a rearrangement of cortical recurrent connectivity towards replayed targets, which is known to suffer from catastrophic forgetting.
\newline\newline
Finally, we found the model suitable to reproduce semantic memory deficits linked to lesioning of the episodic memory system. In particular, we aimed at confirming the standard anterograde semantic amnesia (lesioning MTL before consolidation impairs semantic recall) and, more importantly, retrograde  semantic amnesia \shortcite{duff_2020}. Retrograde semantic amnesia is harder to explain with classic models of semantic consolidation, because the role of MTL is only to hash cortical representations. We first performed a standard anterograde lesion test: a new association was encoded, after which MTL recurrent and replay dynamics were disrupted before consolidation. Under this manipulation, later cortical completion of the associated concept was markedly reduced relative to sham networks (Fig. \ref{fig:behavioural}c).
\newline\newline
However, this standard account does not predict a second effect. Then, we tested whether MTL-semantic also contributes to the retrieval of already consolidated semantic structure. To do so, we allowed a new association to be encoded and consolidated, and only then lesioned the semantic MTL pathway. We then probed cortical recall from partial cues while progressively increasing cortical noise. Sham networks could still recover the semantically related representation, whereas lesioned networks showed a much steeper loss of semantically related cortical activity (Fig. \ref{fig:behavioural}f). Thus, while the model matches the classical role of MTL in new learning (Fig. \ref{fig:behavioural}e), it also captures an effect that standard one-way consolidation models fail to explain: semantic representations in MTL can continue to support semantic recall even after cortical learning has occurred (Fig. \ref{fig:behavioural}f).
\newline\newline
Together, these simulations show that the same structured overlap that improves episodic storage also reproduces behavioural effects that are difficult to explain with standard consolidation models. First, the incorporation of semantic representations into episodic memory provides a natural explanation for the well-established advantage of semantically interpretable stimuli during recall, despite the increased overlap such representations introduce between memories \shortcite{murray_1976, lin_2021}. Second, the replay-based dictionary learning as a mechanism for consolidation predicts that learning experiences sharing a common latent component should accelerate concept formation, reproducing the behavioral benefits of blocked over interleaved learning \shortcite{flesch_2018}. Finally, the bidirectional interaction between episodic and semantic memory systems, implemented through semantic replay, explains why medial temporal lobe damage can impair not only the acquisition of new semantic knowledge but also the retrieval of previously consolidated semantic information \shortcite{duff_2020}.
\begin{figure}
    \centering
    \includegraphics[width=0.95\linewidth]{Figures/Figure_7.png}
    \caption{
Figure 7: Semantic representations in episodic memory explain a range of behavioral phenomena.
\textbf{a}: Diagram of experimentally observed series of words recall, depending on how meaningful the words in the list are (Left). Diagram of experimentally observed memorability of images depending on whether they are shown as original (\textit{intact}) or phase-scrambled (\textit{scrambled}), which erases semantic interpretability of the image (Right).
\textbf{c}: Blocked versus interleaved pre-training. Episodes are presented either in semantically coherent blocks or in an interleaved order while preserving the overall training statistics.
\textbf{d}: Effects of presentation schedule on behaviouraly performance in humans participatns (left) and in the model (right).
\textbf{e}: Anterograde lesion experiment. After pre-training, a novel association is learned and medial temporal replay mechanisms are disrupted before consolidation. Subsequent cortical recall of the learned association is impaired, reproducing the classical dependence of semantic learning on MTL-mediated consolidation.
\textbf{f}: Retrograde lesion experiment. Semantic statistical associations are learned during pre-training, then MTL is lesioned. Cortical semantic recall is then tested from noisy partial cues.
    }
    \label{fig:behavioural}
\end{figure}
\newpage
\section*{Discussion}
We have developed a model of episodic-semantic memory interactions based on two complementary principles. First, episodic replay acts as a form of dictionary learning, enabling neocortex to discover reusable semantic representations from episodic experience. Second, these abstract representations can subsequently be transferred back into the episodic memory system through cortical-driven spontaneous activity, allowing semantic structure to become part of the episodic code. This reciprocal interaction reproduces key experimental observations, including X. Furthermore, the model provides mechanistic insights into the normative role of semantic representations in episodic memory, suggesting that they act as an error-correcting scaffold that improves recall and stabilizes replay dynamics as long as new episodes stay semantically within distribution. This offers a potential explanation for a long-standing tension between theory and experiment: the semantic representations found in the medial temporal lobe should impair attractor dynamics. In summary, we propose a unified framework in which episodic and semantic memory continuously shape one another, each adapting their code to represent the regularities of the other.
\newline\newline
A first contribution of this work is to recast semantic consolidation as a form of dictionary learning. Classical consolidation theories propose that replay enables the gradual emergence of semantic knowledge from episodic experience \shortcite{mcclelland_1995, oreilly_2014}, with many computational implementations relying on supervised or error-driven learning mechanisms to achieve this transformation \shortcite{saxe_2019, singh_2022, sun_2023, spens_2024}. In contrast, our framework suggests that replay may provide candidate structure from which neocortex discovers a set of reusable representations capable of explaining recurring patterns across experiences. This is consistent with increasing evidence indicates that replay is selective, reconstructive, and generative rather than a faithful reinstantiation of prior activity \shortcite{swanson_2020, diekmann_2023}. Furthermore, in the model, semantic representations can emerge without requiring predefined targets, explicit input-output distinctions, or exact reinstatement of previous experiences. More broadly, this perspective provides a unified account of two complementary functions often attributed to systems consolidation: the long-term stabilization of individual episodic memories and the extraction of abstract knowledge across experiences.  Existing computational work has largely treated these objectives separately, with high-level machine learning models focusing on semantic abstraction and generalization, and biologically grounded circuit models primarily addressing the stabilization of individual memories, often inspired by contextual fear conditioning paradigms. How the same consolidation mechanism can support both simultaneously is unknown. Our results suggest that these are different outcomes of the same replay-driven learning process, depending on the episodic complexity of the word (or \textit{atom}) formed. This same process could account for personal semantics \shortcite{renoult_2012}, when the concept formed corresponds to an intermediate complexity level (e.g. the concept of someone's workspace). From this perspective, the level of abstraction of the replayed hippocampal trace itself allows the formation of concepts of different levels of abstraction, supporting a continuum from purely episodic to highly abstract.
\newline\newline
\begin{comment}
First, our results suggest a reinterpretation of replay-based consolidation as a form of dictionary learning. Classical consolidation theories propose that replay enables the gradual transfer of knowledge from episodic to semantic memory \shortcite{mcclelland_1995, oreilly_2014}, with many computational implementations operationalizing this transfer through supervised or error-driven learning mechanisms \shortcite{saxe_2019, singh_2022, sun_2023, spens_2024}. In contrast, our model proposes that replay might serve a different computational role: rather than providing cortical targets, it generates candidate structure from which neocortex can discover recurring latent components of experience. Through local Hebbian plasticity and competition, these components become semantic representations that can be reused across multiple episodes. This perspective naturally addresses several limitations of previous approaches, including the need for predefined input-target distinctions, exact reinstatement of previous activity patterns, or biologically implausible learning rules. Furthermore, it unifies the mechanisms for consolidation of episodic experience and abstract re-usable concepts. The model shows that different sparsity regimes can lead to the consolidation of concepts of different complexity. Importantly, it is also consistent with a growing body of work suggesting that replay contributes not only to memory preservation but also to abstraction, statistical learning, and the extraction of regularities across experiences \shortcite{schapiro_2017, schapiro_2016, sun_2023, barry_2023}. Likewise, increasing evidence indicates that replay is selective, reconstructive, and generative rather than a faithful reproduction of prior activity \shortcite{swanson_2020, diekmann_2023, zhou_2025}. Viewed through this lens, replay may be understood as exposing latent structure embedded within episodic experience, while cortical learning discovers a compact set of reusable semantic building blocks that explain this structure. Interestingly, this interpretation bears a close conceptual resemblance to dictionary learning and sparse coding frameworks, where complex observations are represented through combinations of latent components \shortcite{olshausen_1996}, as well as recent feature-discovery approaches in machine learning that recover semantically meaningful concepts directly from neural activity patterns \shortcite{bricken_2023}. Together, these findings suggest that the primary function of replay may extend beyond transferring memories between systems to discovering the semantic structure through which future experiences are interpreted.
\newline\newline
\end{comment}
A second central finding of this work is a mechanism for transferring abstract representations from one sub-network to another. In the model, spontaneous activity in neocortex shapes medial temporal representations, increasing their level of abstraction during episodic encoding. The idea that neocortex can transfer semantic representations to subcortical areas has inspired several computational models in the recent years. For example, cortex-initiated generative replay has been proposed to align cortical and hippocampal representations \shortcite{van_2020}. Similarly, this assumption underlies models in which cortical-extracted regularities can be used for episodic compression \shortcite{spens_2024_b, nagy_2025} and by the episodic pattern completion machinery \shortcite{whittington_2020, kang_2023, spens_2024, kang_2024}. All of these proposals have in turn been inspired by experimental findings supporting the presence of semantic representations in the episodic memory circuit\shortcite{quiroga_2005, quiroga_2012, aronov_2017, sun_2020}. Furthermore, converging lines of evidence point at neocortex actively influencing hippocampal replays and subsequent consolidation \shortcite{remondes_2004, rothschild_2017, de_sousa_2026}. However, a  biologically plausible circuit model of this process remains yet to be proposed. While still hypothetical, our results provide mechanistic account for the formation of semantic representations in episodic memory relying only on local Hebbian plasticity. Furthermore, the model's mechanism is consistent with experimental findings in which neocortical recurrent engrams serve as a substrate for further synaptic potentiation of downstream corticofugal connections \shortcite{lee_2023} -in particular from prefrontal cortex (PFC) to basolateral amygdala (BLA)- which we have reproduced here.
\newline\newline
Beyond providing a mechanistic account of how semantic representations emerge and become incorporated into episodic memory, our results suggest a normative role for semantic structure within episodic codes. Classical theories of associative memory have long emphasized the importance of sparse and decorrelated representations for maximizing storage capacity \shortcite{hopfield_1982, tsodyks_1988, treves_1991}. From this perspective, semantic overlap between memories would be expected to introduce interference and impair retrieval. Far from being the case, semantic representations can in fact benefit episodic memory \shortcite{murray_1976, tse_2007, lin_2021}. Our results offer one possible explanation for this apparent tension. By providing abstract representations that capture regularities across experiences, semantic concepts act as an error-correcting scaffold during memory retrieval, allowing partially degraded episodic patterns to be reconstructed more accurately. Importantly, this benefit extends beyond recall itself. Because replay relies on the reactivation of stored episodic representations, improved attractor dynamics also lead to more reliable replay events and an increased signal-to-noise ratio during consolidation. In this view, semantic representations do not merely coexist with episodic memories but actively improve their stability, retrieval, and communication with neocortex. More broadly, these findings suggest that semanticization may represent a computational trade-off in which a modest reduction in memory specificity is compensated by substantial gains in robustness, generalization, and consolidation efficiency assuming the world statistics will keep semantically within distribution.
\newline\newline
Taken together, our results aim to reconcile the generally assumed one-way consolidation of circuit models of systems consolidation \shortcite{oreilly_2014,kumaran_2016} with a rich body of experiments demonstrating that semantic knowledge shapes episodic memory. These include reconstructive and schema-dependent recall \shortcite{bartlett_1932}, false memories for semantically related items in the DRM paradigm \shortcite{deese_1959,roediger_1995}, expertise-dependent memory \shortcite{chase_1973}, and improved recall for semantically congruent episodes \shortcite{zoellner_2023}. Recent computational models have successfully captured several of these phenomena through semantic compression \shortcite{nagy_2020} or semantic completion of partial episodic traces \shortcite{fayyaz_2022}. In contrast, our work proposes an explicit hippocampal-neocortical circuit mechanism based on replay, local plasticity, and bidirectional consolidation, providing a biologically grounded framework in which these semantic effects can emerge. While the specific mechanisms remain to be experimentally validated, we hope this model provides a useful circuit-level scaffold for generating mechanistic hypotheses that bridge psychological phenomena and systems-level learning theories.
\newline\newline
Our framework generates several experimentally testable predictions. In an effort to map our episode generation protocol to standard training paradigms in rodents, we propose a paradigm for compositional cues ($A_i$ visual, $B_j$ odor) - reward associations (Figure \ref{fig:predictions_supp}). Then, we propose one prediction for each of the three key mechanisms studied here. First, if dictionary learning is the substrate for semantic learning, we anticiapte the formation of cells tuned to the different components $A_i$, $B_j$ to increase monotonically with the marginal probability of that concept to appear in any episode (Fig. \ref{fig:phase_a}g). A parallel prediction is that holographic manipulation of the formed abstract cells to reflect the conditional probability (semantic) structure, encoding the \textit{world model} in recurrent connections (Fig. \ref{fig:phase_a}h), which would support our mechanism for slow extraction of predictivenes across learned concepts. Second, if semantic representations are transferred back into episodic memory through semantic replay, disrupting cortical activity during consolidation should reduce the abstractness of medial temporal lobe representations, which could be similarly measured via cross-context decodability, selecitivity, or mutual information with the different latent factors presented to the animal. Finally, the model predicts that memories composed of previously learned concepts should consolidate more rapidly than entirely novel experiences. Behaviourally, this would imply that learning the reward association to a novel visual-odor $(A_i, B_j)$ cue would be faster when each individual cue has previously been learned than when the cues are novel (Fig. \ref{fig:normative}d).
\newline\newline
Several limitations of the present work promising directions for future research. First, the model intentionally abstracts away much of the anatomical and physiological complexity of hippocampal-neocortical interactions. While this simplification allows us to isolate the computational principles underlying semantic consolidation and semantic replay, future work could investigate how these mechanisms interact with the rich circuitry of the medial temporal lobe, including the distinct contributions of entorhinal cortex, hippocampus subregions, or amygdala. Second, our implementation focuses on static representations and does not explicitly address the temporal structure of episodic memories. Extending the framework to sequences, trajectories, and temporally extended experiences may provide a natural link to theories of replay, planning, and predictive memory. Finally, while we study a regime whithout imposed pattern separation in hippocampal code (\shortcite{sun_2023, chandra_2025}), it is likely that memories are involve both structured overlapping representations and pattern separated hashing codes \shortcite{chettih_2024}. The computational mechanisms supporting flexible use of these both during encoding and retrieval are yet to be established. More broadly, from a theory of learning perspective, we hope that the present framework provides a foundation for investigating episodic and semantic memory as cooperative components of systems that can simultaneously generalize and memorize. Indeed, the relevance of semantic and episodic memory in modern machine learning algorithms is becoming increasingly more relevant \shortcite{fountas_2025, dong_2025, pink_2025}.
\newpage
\begin{comment}
\section*{Discussion}
In this study, we introduce a cognitive strategy by which episodic memory is improved through leveraging previously extracted semantic representations. Our model posits that, following an episodic-to-semantic consolidation process, semantic-to-episodic memory transfer reshapes the code of episodic memory itself. We further propose a biologically plausible neural circuit that implements such a cognitive model, successfully capturing the associated learning algorithm. Our circuit model aligns with and explains experimental data, and highlights the behavioural advantages of such a bidirectional learning mechanism.
\newline\newline
First, we demonstrated how episodic replay from hippocampal and para-hippocampal (medial temporal lobe, MTL) regions effectively supports semantic consolidation in cortical (CTX) networks. While this initial step of consolidation corresponds to the most widely studied phase of learning across subcortical-cortical interactions \shortcite{alvarez_1994, mcclelland_1995, oreilly_2014, sun_2023, guerreiro_2024, spens_2024} (to name just a few), our model exhibits key differences from previous models. To start, most models assume a faithful replay of awake subcortical and cortical activity, with abstraction happening directly in CTC$\leftarrow$CTX (neocortex to neocortex) connections. This is problematic for two reasons: first, the assumption that subcortical replay can reinstate perfectly previous cortical activity is not given \shortcite{swanson_2020}. Secondly, this relies on meaningful (to some degree) representations already existing in neocortex, when it is precisely via semantic consolidation that these representations should be created in the first place. For example, in \shortciteA{chrysanthidis_2022}, semantics are also extracted by forming block-like structure of recurrently connected neurons, but the high specificity of these neurons is already given. Furthermore, many studies assume a copy-paste mechanism between cortex and hippocampus, either directly via one-to-one connections \shortcite{alvarez_1994, tome_2022, guerreiro_2024, spens_2024} or perfect pattern-separating random connections \shortcite{sun_2023}, with less emphasis placed on how synapses mapping one to the other can be formed with biologically plausible learning mechanisms (although previous effort has been made in the context of backpropagation-like error-driven learning rules \shortcite{singh_2022}). In contrast, here we have proposed that increasing sparsity levels during replay can lead to CTX$\leftarrow$MTL synapses becoming selective to frequent episodic sub-patterns. This effectively decorrelates the training input provided to these connections, which could otherwise lead to representational collapse using standard Hebbian mechanisms. 
\newline\newline
Next, we proposed a mechanism whereby the semantic structure imprinted in CTX$\leftarrow$CTX connections can be transferred into the coding scheme of MTL. Both the notion that semantics are extracted in CTX \shortcite{mcclelland_1995} and that the MTL encodes a latent representation of the environment \shortcite{quiroga_2005, aronov_2017, behrens_2018} have been present in the literature for a long time, but a mechanistic account of how these representations emerge in MTL remains unknown. We put forward the hypothesis that, after extracted in CTX, spontaneous cortical activity could transfer information encoded in cortical synapses to MTL$\leftarrow$CTX synapses, inducing a semantic code in MTL. We called this process semantic-to-episodic consolidation, and conceptualized it as a mirror of episodic-to-semantic consolidation. As such, instead of transferring episodic information contained in recurrent MTL connections to CTX representations, it transfers semantic information present in recurrent CTX connections to MTL representations. The idea that cortex-to-hippocampus teaching is an underexplored topic in complementary learning systems literature was recently highlighted in \shortciteA{liu_2024}, and implicitly explored in \shortciteA{spens_2024}. In the latter study, the authors explored the possibility of using a Hopfield network to store pairs of sensory-latent representations for episodic recall, in a similar manner that our MTL-sensory and MTL-semantic neurons form joint attractors in MTL recurrent connections. Again, however, the process by which these joint representations were formed was not fully specified, with the model instead employing a copy-paste strategy to generate them. A model that also studied the formation of joint sensory-semantic representations as attractors in hippocampus is the Tolman Eichenbaum Machine \shortcite{whittington_2020}. In this study, the information in the sensory input was already highly semantized (i.e. using a one-hot encoding of sensory features), and the abstract graph position of the stimuli also given a priori. The mechanism for obtaining these representations in the first place was also not explored. In addition to establishing a mechanism for concept cell formation in MTL, our results show that this form of representations is more suitable for rapid learning than mixed representations, both in In-Distribution and Compositional Out-Of-Distribution environmental statistics. This allows MTL to compositionally encode never-seen episodes as long as these correspond to rearrangements of previously seen semantic representations. This is in agreement to previous results in continual learning in artificial neural networks \shortcite{chaudhry_2020}. We also showed how maintaining sensory low-level representations of the environment is crucial for encoding episodes that are Semantic Out-Of-Distribution (that is, when new episodes cannot be encoded using previous semantic representations). 
\newline\newline
Our model has shown how even in non-trivial, highly overlapping scenarios like a random projection of the original input one can obtain concept cells in MTL. Moreover, the model offers a mechanism via which the formation of these representations is favored by blocked (as opposed to interleaved) learning, as demonstrated in humans \shortcite{flesch_2018}. In our model, blocked learning provides the MTL with many different examples of episodes containing the same concept, which are averaged out in its recurrent connections, facilitating the replay of its corresponding prototypical representation \shortcite{posner_1968}. The model also explains why recall of episodes with a higher semantic component is better than that of episodes that, from the perspective of the learned semantic structure, appear random \shortcite{lin_2021}. In classic models of recall based on attractor models, overlaps between episodes is seen as detrimental, reducing the network capacity \shortcite{hopfield_1982, treves_1994}. However, semantic representations have a structured overlap, where the overlap is not stochastic but instead corresponds to the two episodes containing semantic overlaps. Thus, rather than reducing recall, this overlap acts as an error-correcting code \shortcite{albesagonzalez_2025}. Moreover, semantic representations not only show better recall, but also actively improve the recall of the sensory component. Similarly, semantic representations induce replaying activity patterns that resemble more activity during wakefulness. These results can help explain why learning schema-congruent information is much faster than that of completely new episodes \shortcite{tse_2007}. While many models have focused on explaining how faster neocortical learning could happen at the cortex level \shortcite{mcclelland_2013, guerreiro_2024}, our results suggest that this could be in part because the replayed activity in the first place is less noisy. Understanding the mechanisms of positive transfer learning is one of the cornerstones of modern artificial intelligence, where inducing models to use previously learned information to learn new tasks is still a challenge \shortcite{iman_2023}.
\newline\newline
Following this potential link between replay involving semantic representations and continual learning, we showed that closing the loop (performing episodic-to-semantic consolidation of episodes that contain both sensory and semantic representations) allows the formation of higher-order compositional representations in CTX \shortcite{kurth_2023}, in turn inducing a complete semantic description of the causal links across concepts in different levels of abstraction \shortcite{barlow_1972}. Not only higher-order representations (for example pizza is a conjunction of the different ingredients it contains) were developed, but CTX$\leftarrow$CTX connections captured the semantic relationships across these different hierarchical levels (such as the fact that pizza is much more related to cheese than cheese is to pizza). Crucially, these forms of asymmetric relationships cannot be captured by classic models of Hebbian learning, as recently shown in \shortciteA{albesagonzalez_2025}. Finally, our framework might be related to the phenomenon of infantile amnesia \shortcite{alberini_2017}. In the model, the semantic load quantifies how many distinct concepts are bound within a neural representation, ranging from low (highly conceptual) to high (highly episodic) values. Furthermore, we showed that consolidating representations with a high semantic load (those combining many concepts) requires pre-existing semantic structure to scaffold them. It could be that early in development, when semantic representations are still forming, this scaffolding is absent, making the consolidation of rich, episodic memories difficult. Consequently, infants may initially consolidate only simple, concept-like representations (e.g., family members or locations), while conjunctive memories combining these elements fail to stabilize. As semantic representations mature, episodes involving multiple known concepts can finally be consolidated. It would be interesting to study whether this coincides with the end of infantile amnesia.
\newline\newline
The model faces several limitations at different levels. For example, the model omits hash-like, episode-specific representations in episodic memory (as have been observed experimentally \shortcite{kolibius_2023, chettih_2024} and included in previous models \shortcite{sun_2023, chandra_2025}). An open question in neuroscience is how episodic representations solve the trade-off between perfect pattern separation and using overlapping patterns, more aligned with the hippocampus as a cognitive map that includes behaviourally-relevant latent representations. Further studies using a combination of shared and non-shared neuronal subspaces across similar episodes could explore a combination of both strategies might facilitate learning and optimal behaviour (efforts in this direction have been made before \shortcite{schapiro_2017}). Additionally, while we have slightly moved away from completely orthogonal input representations usually used in biologically-plausible mechanistic models of consolidation, the input used here was still relatively simple, with more complex representations still requiring gradient-based methods. How the brain might use local plasticity rules to learn from more realistic input remains an open problem.
\newline\newline
The results presented here could finally be of significance for modern artificial neural networks. In the recent years, many authors have highlighted episodic memory as a missing piece in many Artificial Intelligence (AI) paradigms .  Our model suggests that episodic memory should not simply be added as a fast external buffer, but rather be continuously reshaped by semantic knowledge. This bidirectional interaction yields two advantages directly relevant for AI: (i) more robust recall of past experiences through semantic “error correction,” and (ii) the ability to rapidly learn new combinations of previously known concepts, supporting compositional generalization. Such designs may allow modern neural networks to leverage prior structure in order to accelerate continual learning and adapt flexibly to new environments.
\newline\newline
Altogether, our study shows that bidirectional consolidation can be captured by a biologically plausible circuit model relying only on local plasticity rules. We demonstrated how sparsity during replay enables cortical extraction of compositional structure, how semantic feedback supports the emergence of concept cells in MTL, and how the presence of semantic codes enhances recall, replay quality, and higher-order concept learning. These results extend standard consolidation models by providing a mechanistic explanation for the formation of semantic representations in MTL and their role in scaffolding both new episodic and semantic memories.
\end{comment}
\newpage
\section*{Methods}\label{methods}
\subsection*{Episode Generation Protocol}
The Episode Generation Protocol (EGP) is a generative process that samples neuronal activity in the sensory layer. Our proposed EGP follows a similar approach to that described in \shortciteA{albesagonzalez_2025}, but is outlined here again for completion:
\subsection*{Episodes, episode attributes and episode concepts}
All \textit{episodes} share a structure, such that every episode contains one \textit{episode concept} per \textit{episode attribute}. This means that, given episode attributes $A$ and $B$, with each episode attribute representing a set of episode concepts $a\in A, b\in B$, then the set of all possible episodes is defined as:
\begin{equation}
    E = \left\{ \{a, b\} \mid a \in A, b \in B \right\}
\end{equation}
In other words, an episode $e\in E$ is a collection of episode concepts $a$, $b$ such that element $a$ belongs to episode attribute  $A$ and element $b$ belongs to episode attribute $B$. In this study we have fixed $a$ to have 5 possible values $a\in \{A_1, A_2, A_3, A_4, A_5\}$ and likewise for $b\in \{B_1, B_2, B_3, B_4, B_5\}$.
\newline\newline
\subsection*{Episode Probability Distribution}
As a generative process, our EGP samples episodes as a previous step to sampling sensory input. We do this by fixing a probability distribution over episodes:
\begin{equation}
    1 = \sum_{A_i, B_j}P(e = \{A_i, B_j\}) 
\end{equation}
such that $P(e = \{A_i, B_j\}) \equiv P(i,j)$ imposes the likelihood of an episode with latents $(A_i, B_j)$ to be sampled when generating episodes.
\subsection*{Episode-Input Mapping}
Ultimately, the EGP samples sensory layer activity. For this reason, one also has to define how a sampled episode is mapped into sensory input $X^\textrm{SEN}$. Every concept $c$ in $A\cup B$ has a set of associated neurons $\textrm{SEN}_{c}$ in $X^\textrm{SEN}$ such that, given an episode $e$
\begin{equation}
     X^\textrm{SEN}_{i} = \begin{cases} 
      1 & \textrm{if} \;\; i \in \textrm{SEN}_{c} \;\;\textrm{and}\;\; c \in e\\ 
      0 & \textrm{else} \\
   \end{cases}
\end{equation}
In order to define $\textrm{SEN}_c$, we specify the number of attributes $N_\textrm{attributes}$, in this case 2 ($A$ and $B$). Then, we define the number of concepts per attribute ($|A|$ and $|B|$, here 5 and 5, from this follows the fact that $A_{i \in [1, 5]}$ and $B_{j \in [1, 5]}$). Then, we fix the number of neurons per concept $N_\textrm{concept}^A$ and $N_\textrm{concept}^B$, and define $\textrm{SEN}_{c}$ as:
\begin{equation}
    \textrm{SEN}_{c} = \{i : \textrm{start}(c) \leq i < \textrm{stop}(c)\}
\end{equation}
with $\textrm{start}(c) = (i - 1) \cdot N_\textrm{concept}^A$ if $c = A_i \in A$ and $\textrm{start}(c) =|A|\cdot N_\textrm{concept}^A + (j - 1) \cdot N_\textrm{concept}^B$ if $c = B_j \in B$, and similarly $\textrm{stop}(c) = i \cdot N_\textrm{concept}^A$ if $c = A_i \in A$ and $\textrm{stop}(c) =|A|\cdot N_\textrm{concept}^A + j \cdot N_\textrm{concept}^B$ if $c = B_j \in B$ (i.e. we generate a list of one-hot encodings for each concept $c$).
\newline\newline
To account for variability between different presentations of the same episode (not all episodes, even though they contain the same concepts, will be exactly the same), we further add some stochasticity by randomly picking $N_\textrm{swap}$ inactive neurons and $N_\textrm{swap}$ active neurons and inverting their activity (a total of $2N_\textrm{swap}$ neurons randomly change their activity). This ensures that activity sparsity in the sensory layer is maintained (the number of swaps from 0 to 1 is the same as the number of swaps from 1 to 0).
\subsubsection*{Input temporal structure}
Generated input $X_\textrm{input}$ has dimensions $(N_\textrm{days}, T_\textrm{day}, N_\textrm{SEN})$. Every day, a total of number $T_\textrm{day}/T_\textrm{episode}$ episodes (of $T_\textrm{episode}$ length) are sampled, and within each timestep an independent noise realization is sampled. 
\subsubsection*{Semantic Structure of an EGP}
Following \shortciteA{albesagonzalez_2025}, we refer to  \textit{semantic field theory} \shortcite{bussmann_2006} in order to define what is the \textit{semantic structure} of an EGP. According to this theory, the meaning of a word is not isolated but dependent on its relation to the rest of the words. While our task is not one of language, we can use this same paradigm to define the \textit{meaning} of episode concepts. In this sense, the meaning (semantics) of our episode concepts depends on how they are related to the rest (also note this is how semantics are established in a language dictionary, where each concept is defined via its relationship with the rest of concepts). 
\newline\newline
Following this proposal, we use the conditional probabilities of being present in an episode between episode concepts as a proxy for the semantic structure of an EGP. In other words, extracting the semantics of an EGP is equivalent to extracting how likely is one episode concept $y$ to be present in an episode $e$ if an episode concept $x$ is also present. We define the Semantic Structure $\mathcal{SS}$ of an EGP with $N^\textrm{concepts}$ different concepts $c_i$, as as a $N^\textrm{concepts}\times N^\textrm{concepts}$ matrix of the form:
\begin{equation}
    \mathcal{SS}_{ij} = P(c_i \in e|c_j \in e) 
    \label{eq:semantic_structure}
\end{equation}
For example, the intuition that pizza has a stronger semantic connection to cheese than cheese to pizza is reflected in their conditional probabilities (it is very common to have cheese in a pizza, but much less to have pizza in any dish containing cheese).  
\subsection*{Circuit Model}
\subsubsection*{Architecture and nomenclature}
\paragraph*{Regions and neural activity}
The circuit model contains 3 main regions: sensory cortex (SEN), neocortex (CTX), and medial temporal lobe (MTL). In turn, MTL is split into two sub-regions: MTL-sensory and MTL-semantic. We use $X$ to denote the neural activity, which is written as $X^\textrm{region}$ (for example $X^\textrm{CTX}$) when the region is to be made explicit. The same follows for pre-activation (synaptic) input, which is denoted by $\hat{X}$ (or $\hat{X}^\textrm{region})$. The activity of neuron $i$ in a region is $X^\textrm{region}_i$. The number of neurons in a region is denoted by $N_\textrm{region}$ (e.g. $N_\textrm{CTX}$).
\paragraph{Connectivity}
Connectivity matrices connecting a region \textrm{pre} to a region \textrm{post} take the form $W^{\textrm{post}\leftarrow\textrm{pre}}$, and a synapse connecting neuron $j$ in region pre with neuron $i$ in region post is $W^{\textrm{post}\leftarrow\textrm{pre}}_{ij}$. For example, $W^{\textrm{CTX}\leftarrow\textrm{MTL}}_{13}$ represents the connection from neuron 3 in MTL to neuron 1 in CTX. Connectivity is classified as either feed-forward (connecting two different regions), which includes  $W^{\textrm{MTL-sensory}\leftarrow\textrm{SEN}}$, $W^{\textrm{CTX}\leftarrow\textrm{MTL}}$, and $W^{\textrm{MTL-semantic}\leftarrow\textrm{CTX}}$, or recurrent (connecting a region with itself), which includes $W^{\textrm{CTX}\leftarrow\textrm{CTX}}$, and $W^{\textrm{MTL}\leftarrow\textrm{MTL}}$. All connections are plastic except $W^{\textrm{MTL-sensory}\leftarrow\textrm{SEN}}$, which is fixed. Given a connectivity matrix $W^{\textrm{post}\leftarrow\textrm{pre}}$, we call receptive field of a postsynaptic neuron $i$ $\textrm{RF}^{\textrm{post}\leftarrow\textrm{pre}}_i$ the vector representing the connections from region pre to postsynaptic neuron $i$
\begin{equation}
    \textrm{RF}^{\textrm{post}\leftarrow\textrm{pre}}_i = \{{W^{\textrm{post}\leftarrow\textrm{pre}}}_{kj}\} \;\; : \;k = i 
\end{equation}
\paragraph{Synaptic Initialization}
\paragraph{Sensory cortex to MTL-sensory}
$W^{\textrm{MTL-sensory}\leftarrow\textrm{SEN}}$ is initialized as the identity, mapping to MTL-sensory an exact copy of $X^\textrm{SEN}$.
\paragraph{Feed-Forward connections}
Feed-forward connections are sampled from a Gaussian distribution with mean 0 and standard deviation $W^{\textrm{post}\leftarrow\textrm{pre}}_\sigma$ as a network parameter.
\begin{equation}
    W^{\textrm{post}\leftarrow\textrm{pre}}_{ij} \sim \mathcal{N}(0, W^{\textrm{post}\leftarrow\textrm{pre}}_\sigma)
\end{equation}
\paragraph{Recurrent connections}
Recurrent connections are initialized at 0. MTL connections are reset every \textit{day} (see Algorithm \ref{alg:train})
\paragraph{Mature and immature neurons} Neurons in the model can be either \textit{mature} or \textit{immature}. The intuition is: do these neurons contain meaningful representations via input received from another region (mature) or are mostly non-selective and random (immature)? More details on the dynamics and plasticity of mature and immature neurons can be seen below, but to summarize: (i) immature neurons have higher excitability during feed-forward input processing (not during pattern completion) and (ii) immature neurons have a higher learning rate in their incoming feed-forward connections and its recurrent connections frozen. Because of their higher excitability, during receptive field formation in feed-forward synapses, immature neurons win the competition for postsynaptic activity only if the presynaptic pattern is too far from the receptive fields that have already formed. By doing this, if a presynaptic pattern is similar to previously formed representations, it slightly changes the already formed representations. If not, a set of immature neurons win the competition, develop receptive field due to higher plasticity, and then become mature. This allows one-shot formation of new postsynaptic representations of presynaptic input when an input pattern is too different from those presented to the network before, and then refining these representations with a slower learning rate. Recurrent connectivity in CTX, which reflects statistical regularities across the representations formed, is only developed between mature neurons that actually code for an environmental variable, to avoid recurrent connections representing spurious correlations. Maturity in a region is denoted by $\textrm{IM}^\textrm{region}$, where $\textrm{IM}^\textrm{region}_i = 1$ if neuron $i$ in the region is immature, and 0 otherwise.
\subsubsection*{Outgoing and incoming homeostasis across types of connections}
Below, we formalize two regimes of homeostatic plasticity: outgoing homeostasis and incoming homeostasis. Outgoing homeostasis guarantees that the sum of the outgoing connections from neuron $j$ in a connectivity matrix $W$ does not exceed a certain value $W_\textrm{max-out}$ ($\sum_iW_{ij} \leq W_\textrm{max-out} \forall j$). Similarly, incoming homeostasis guarantees a bound in incoming connections ($\sum_jW_{ij} \leq W_\textrm{max-in} \forall i$). Depending on the type of connectivity we choose:
\begin{itemize}
    \item outgoing homeostasis for CTX$\leftarrow$CTX: this makes connectivity matrix reflect semantic structure as defined in \eqref{eq:semantic_structure}, as connection from neuron coding from concept $j$ to neuron coding for concept $i$ becomes proportional to the conditional probability of neuron $i$ firing given neuron $j$ fires \shortcite{albesagonzalez_2025}
    \item incoming homeostasis in feed-forward connectivity: in this homeostatic regime, connection from neuron $j$ to neuron $i$ becomes proportional to $p(j \; \textit{active} | i\; \textit{active})$ \shortcite{albesagonzalez_2025}. The \textit{trick} is that if an initial one-shot receptive field makes neuron $i$ selective to concept $c$, receptive fields approximate the average presynaptic activity pattern (prototype) given the concept c is present (neuron $i$ is active), thus reinforcing higher selectivity in noisier regimes.
    \item Homeostatic reset in MTL$\leftarrow$MTL connections. In order to account for symmetric plasticity kernels in hippocampus \shortcite{mishra_2016, li_2023}, we simply have Hebbian learning in MTL$\leftarrow$MTL. Homeostasis is obtained by resetting MTL recurrent connections to 0 after each sleep cycle.
\end{itemize}
\subsection*{Network Operations}
\subsubsection*{Neuronal activation}
Synaptic input is mapped to neural activity via a \textit{K-winners-take-all} mechanism, in which the $K$ neurons with the highest pre-activation input $\hat{X}$ set their activity to 1, and the rest to 0. This can be written as
\begin{equation}
    X = \textrm{top-}K\big(\hat{X}\big)
\end{equation}
In  order to determine $K$ in a particular region, first we define
\begin{equation}
    K^\textrm{region}_\textrm{concept} \leq K^\textrm{region}_\textrm{episode}
\end{equation}
 such that each concept is represented by $K^\textrm{region}_\textrm{concept}$ neurons, and a typical episode during wakefulness contains $K^\textrm{region}_\textrm{episode}$ active neurons. In this context, given $K$ active neurons in a region, we define the semantic load $S_L$ to be the number of concepts represented by $K$ neurons, such that
 \begin{equation}
     K = S_L \cdot K^\textrm{region}_\textrm{concept}
     \label{eq:semantic_load}
 \end{equation}
In other words, by fixing the semantic load, one fixes how many neurons in total are active, making $K$ higher if more concepts are being represented. The maximum number of concepts that can be encoded in an episode is thus:
 \begin{equation}
     N^\textrm{region}_\textrm{concept} = \frac{K^\textrm{region}_\textrm{episode}}{K^\textrm{region}_\textrm{concept}}
 \end{equation}
 Therefore, we can define \textsc{Activate} (Algorithm \ref{alg:activate}) as an algorithm that maps presynaptic input $\hat{X}$, given a semantic load $S_L$, into neural activation $X$ in the following way:
\begin{algorithm}[h!]
\caption{Neuronal Activation}\label{alg:activate}
\begin{algorithmic}[1]
\Function{Activate}{$\hat{X}$, $\textrm{region}$, $S_L$}
    \State Compute number of active neurons:
        \[
            K \gets S_L \times K^\textrm{region}_\textrm{concept}
        \]
    \State Select top-$K$ entries:
        \[
            X \gets \textrm{top-}K(\hat{X})
        \]
    \State \Return $X$
\EndFunction
\end{algorithmic}
\end{algorithm}

\subsubsection*{Feed-forward input processing}
Input processing is computed for a single pre and post region at a time, with pre-activation input $\hat{X}^\textrm{post}$ obtained via
\begin{equation}
    \hat{X}^\textrm{post} = W^{\textrm{post}\leftarrow\textrm{pre}} \cdot X^\textrm{pre} + b^\textrm{post}\textrm{IM}^\textrm{post}
\end{equation}
$b^\textrm{post}$ is the extra excitability that immature neurons have in region post.
\subsubsection*{Pattern completion (recurrent input processing)}
Given a recurrent network, pattern completion is typically defined as an iterative update on neural activity $X^\textrm{region}$ that is initiated in an initial state $X^\textrm{region}_0$:
\begin{equation}
    X^{\textrm{region}}_{k+1} = \textrm{top-}K(W^{\textrm{region}\leftarrow\textrm{region}}\cdot X^\textrm{region}_{k})
\end{equation}
with $W^{\textrm{region}\leftarrow\textrm{region}}$ the region recurrent connectivity matrix, $K$ dependent on the specified semantic load $S_L$, and $k$ a timestep operating on a smaller timescale,  than the overall timescale $t$. This means that within a timestep $t$, there are pattern\_complete\_iterations iterations on temporal index $k$. This is used to define \textsc{PatternComplete}  as Algorithm \ref{alg:pattern_complete}:
\begin{algorithm}[h!]
\caption{Pattern Completion}\label{alg:pattern_complete}
\begin{algorithmic}[1]
\Function{PatternComplete}{$X_0$, $S_L$,   \texttt{pattern\_complete\_iterations}}
    \State $X \gets X_0$
    \For{$\text{iteration} = 1$ to $\texttt{pattern\_complete\_iterations}$}
        \State $\hat{X} \gets W^{\textrm{region} \leftarrow \textrm{region}} \cdot X$
    \State $X \gets \textsc{Activate}\bigl(\hat{X}, \textrm{region}, S_L \bigr)$
    \EndFor
    \State \Return $X$
\EndFunction
\end{algorithmic}
\end{algorithm}
\subsubsection*{Hebbian Learning}
The connection between a presynaptic neuron $j$ and a postsynaptic neuron $i$ subject to Hebbian learning is updated as:
\begin{equation}
            \Delta W_{ij} = \lambda_{ij}^{\textrm{post} \leftarrow \textrm{pre}}
            X_i^\textrm{post} X_j^\textrm{pre}
\end{equation}
for a presynaptic region pre and a postsynaptic region post, where $\lambda_{ij}^{\textrm{post} \leftarrow \textrm{pre}}$ is the learning rate from neuron $j$ in region pre to neuron $i$ in region post.
\paragraph*{Maturity-dependent learning rate}
In order to combine one-shot and statistical learning (see \shortciteA{albesagonzalez_2025}), the learning rate of connections depends on the \textit{maturity} of neurons, which is given by $\textrm{IM}$. Neurons start in an \textit{immature} state $\textrm{IM}_i = 1$, allowing them to form initial \textit{connectivity guessings}. Once this initial connectivity is formed, the learning rate is reduced to capture the statistical regularities of the neural activity patterns (the neuron becomes \textit{mature}, $\textrm{IM}_i = 0$). To model this, we define a quick learning rate
\begin{equation}
        \lambda^{\textrm{post}\leftarrow\textrm{pre}}_\textrm{quick, pre/post} = \frac{W^{\textrm{post}\leftarrow\textrm{pre}}_\textrm{max-in/max-out}}{K^\textrm{pre/post}} 
        \label{eq:quick_lambda}
\end{equation}
where again $W^{\textrm{post}\leftarrow\textrm{pre}}_\textrm{max-in}$ ($W^{\textrm{post}\leftarrow\textrm{pre}}_\textrm{max-out}$) or is a network parameter determining the maximum sum of incoming (or outgoing) connections the matrix $W^{\textrm{post}\leftarrow\textrm{pre}}$ has for every postsynaptic (or presynaptic) neuron. Then, at each Hebbian update, we have for incoming-bounded connections:
\begin{equation}
    \lambda_{ij}^{\textrm{post} \leftarrow \textrm{pre}} = (1 - \textrm{IM}^\textrm{post}_i)\lambda^{\textrm{post} \leftarrow \textrm{pre}}_\textrm{slow} + \textrm{IM}^\textrm{post}_i\lambda^{\textrm{post}\leftarrow\textrm{pre}}_\textrm{quick, post}
\end{equation}
and for outgoing-bounded connections:
\begin{equation}
    \lambda_{ij}^{\textrm{post} \leftarrow \textrm{pre}} = (1 - \textrm{IM}^\textrm{post}_j)\lambda^{\textrm{post} \leftarrow \textrm{pre}}_\textrm{slow} + \textrm{IM}^\textrm{post}_j\lambda^{\textrm{post}\leftarrow\textrm{pre}}_\textrm{quick, pre}
\end{equation}
where $\lambda^{\textrm{post} \leftarrow \textrm{pre}}_\textrm{slow}$ is a network parameter indicating the slow learning rate (used throughout all simulations except for the first plasticity event in which the neuron transitions from immature to mature).
\newline\newline
Note how the quick learning rate (Eq. \eqref{eq:quick_lambda}) is defined such that when the neuron transitions from immature to mature, with the  \textit{first shot of learning}, the condition $W^{\textrm{post}\leftarrow\textrm{pre}}_\textrm{max-in}$, for incoming normalized weights, is met:
\begin{equation}
    \sum_j W_{ij} = \sum_j W_{ij}(t = 0) + \sum_j \Delta W_{ij} \approx \;\; 0 (W_{ij} \;\; \textrm{Gaussian})  +\lambda^{\textrm{post}\leftarrow\textrm{pre}}_\textrm{quick, post}\cdot K^\textrm{post} = W^{\textrm{post}\leftarrow\textrm{pre}}_\textrm{max-in}
\end{equation}
(and similarly for outgoing normalized weights). Also note how for incoming-bounded connections, as are feed-forward synapses, this makes the formed receptive field proportional to $X\textrm{pre}$, (making it maximally selective to it). For outgoing-bounded connections, here recurrent cortical synapses, this is equivalent to making $X^\textrm{post}$ the initial guess in meaning of neuron $j$.
\paragraph*{Hebbian learning algorithm}
Altogether, this defines \textsc{Hebbian} (Algorithm \ref{alg:hebbian}) as the following update in $W^{\textrm{post} \leftarrow \textrm{pre}}$:
\begin{algorithm}
\caption{Hebbian Learning}\label{alg:hebbian}
\begin{algorithmic}[1]
\Function{Hebbian}{$\textrm{post}$, $\textrm{pre}$}
    \State Obtain learning rate matrix:
    \If{$\textrm{post} \neq \textrm{pre}$}
        \[
            \lambda_{ij}^{\textrm{post} \leftarrow \textrm{pre}}
            \gets
            (1-\textrm{IM}^\textrm{post}_i)\,\lambda^{\textrm{post} \leftarrow \textrm{pre}}_\textrm{slow}
            + \textrm{IM}^\textrm{post}_i\,\lambda^{\textrm{post}\leftarrow\textrm{pre}}_{\textrm{quick, post}}
        \]
    \ElsIf{$\textrm{post} = \textrm{pre} = \mathrm{CTX}$ AND sleep}
        \[
            \lambda_{ij}^{\textrm{post} \leftarrow \textrm{pre}}
            \gets
            (1-\textrm{IM}^\textrm{post}_j)\,\lambda^{\textrm{post} \leftarrow \textrm{pre}}_\textrm{slow}
            + \textrm{IM}^\textrm{post}_j\,\lambda^{\textrm{post}\leftarrow\textrm{pre}}_{\textrm{quick, pre}}
        \]
    \Else
        \[
            \lambda_{ij}^{\textrm{post} \leftarrow \textrm{pre}} 
            \gets 
            \lambda^{\textrm{post}\leftarrow\textrm{pre}}_0
        \]
    \EndIf
    \State Compute Hebbian update with activity outer product:
        \[
            \Delta W \gets \lambda^{\textrm{post} \leftarrow \textrm{pre}} \odot 
            \Bigl(X^{\textrm{post}} \otimes X^{\textrm{pre}}\Bigr)
        \]
    \State Update weights:
        \[
            W^{\textrm{pre} \leftarrow \textrm{post}} 
            \gets 
            W^{\textrm{pre} \leftarrow \textrm{post}} + \Delta W
        \]
\EndFunction
\end{algorithmic}
\end{algorithm}
\subsubsection*{Homeostatic Plasticity}
We use multiplicative normalization that guarantees, for every weight matrix $W^{\textrm{post}\leftarrow\textrm{pre}}$, the total sum of incoming or outgoing connections is capped to a certain value. Following \shortciteA{albesagonzalez_2025}, we use \textit{incoming homeostasis} as:
\begin{equation}
W_{ij}^{k+1}=
\begin{cases}
 W_{ij}^{k}\cdot\frac{W_\textrm{max-in}}{\sum_{j}W^{k}_{ij}} & \text{if }\sum_{j}W^{k}_{ij} = W_\textrm{max}^\textrm{in} + \epsilon_i, \;\; \epsilon_i > 0 \\[1.2ex]
W^k_{ij}& \textrm{else}
\end{cases}
\end{equation}
where $W^k_{ij}$ denotes the value of synapse $ij$ before applying homeostasis and $W^{k+1}_{ij}$ the value after applying homeostasis (for readability, we simply write $W \equiv W^{\textrm{post}\leftarrow\textrm{pre}}$).
Similarly, we have \textit{outgoing homeostasis} to be defined by:
\begin{equation}
W_{ij}^{k+1} =
\begin{cases}
  W_{ij}^{k}\frac{W_\textrm{max-out}}{\sum_{i}W^k_{ij}} & \text{if }\sum_{j}W^k_{ij} = W_\textrm{max}^\textrm{out} + \epsilon_j, \;\; \epsilon_j > 0 \\[1.2ex]
W^k_{ij} & \textrm{else}
\end{cases}
\end{equation}
This allows, for a presynaptic region pre and a postsynaptic region post, defining \textsc{Homeostasis} as Algorithm \ref{alg:homeostasis}.
\subsubsection*{Fixed points of synaptic dynamics}
The combination of \textit{K-winners-take-all} ($\textrm{top-}K$) activation function with Hebbian and homeostatic plasticity as defined here has been recently shown \shortcite{albesagonzalez_2025} to make synaptic connections converge to conditional firing probabilities. For outgoing homeostasis, the fixed points are:
  \begin{equation}
     W_{ij}^{\textrm{post}\leftarrow\textrm{pre}, \; \textrm{out}}(\infty) \equiv \frac{W_\textrm{max-out}^{\textrm{post}\leftarrow\textrm{pre}}}{K^\textrm{post}}p(X^\textrm{post}_i = 1| X^\textrm{pre}_j = 1)
 \end{equation}
 where $K^\textrm{post}$ corresponds to the number of active neurons during plasticity. Given that the conditional probability can be at most 1, we define in this scenario $W_\textrm{max}$ as $\frac{W_\textrm{max-out}^{\textrm{post}\leftarrow\textrm{pre}}}{K^\textrm{post}}$. Similarly, for incoming homeostasis, the fixed points are
   \begin{equation}
   W_{ij}^{\textrm{post}\leftarrow\textrm{pre}, \; \textrm{out}}(\infty) \equiv \frac{W_\textrm{max-in}^{\textrm{pre}\leftarrow\textrm{post}}}{K^\textrm{pre}}p(X^\textrm{pre}_j = 1| X^\textrm{post}_i = 1)
 \end{equation}
 with $W_\textrm{max}=\frac{W_\textrm{max-in}^{\textrm{pre}\leftarrow\textrm{post}}}{K^\textrm{pre}}$. These results have an important associated intuition, which is that if neurons code for individual concept (have a concept selectivity close to 1), synaptic connections reflect the associated Semantic Structure \shortcite{albesagonzalez_2025} under outgoing homeostasis (as happens in CTX). In MTL-semantic$\leftarrow$CTX and CTX$\leftarrow$MTL, which follow incoming homeostasis, when initial one-shot receptive fields induce a selectivity to a concept close to 1, feed-forward synapses slowly approximate to the associated prototype, where each synapse is becomes proportional to the contribution of each presynaptic neuron to the concept encoded in the postsynaptic neuron.
\begin{algorithm}
\caption{Homeostasis}\label{alg:homeostasis}
\begin{algorithmic}[1]
\Function{Homeostasis}{$\textrm{post}$, $\textrm{pre}$}
\State $W \gets W^{\textrm{post}\leftarrow\textrm{pre}}$
    \For{each post-synaptic unit $i$}
        \If{$\sum_j W[i,j] > W^{\textrm{post}\leftarrow\textrm{pre}}_\textrm{max-post}$}
            \State Normalize:
                \[
                W[i,:] \gets W[i,:] 
                \times
                \frac{W^\textrm{}}{\sum_j W[i,j]}
                \]
        \EndIf
    \EndFor
    \For{each pre-synaptic unit $j$}
        \If{$\sum_i W[i,j] > \text{max\_pre\_connectivity}$}
            \State Normalize:
                \[
                W[:,j] \gets W[:,j] 
                \times
                \frac{\text{max\_pre\_connectivity}}{\sum_i W[i,j]}
                \]
        \EndIf
    \EndFor
\EndFunction
\end{algorithmic}
\end{algorithm}
\subsubsection*{Replay}
We define \textrm{Replay} as a learning mechanism involving a pre and a post region, such that $W^{\textrm{post}\leftarrow\textrm{pre}}$ extract statistical structure from $W^{\textrm{pre}\leftarrow\textrm{pre}}$ connections. A replay operation takes as argument a specific semantic load $S_L$, which determines the conceptual complexity of the extracted regularities. To do this, we first pattern complete an initial presynaptic pattern $X^\textrm{pre}_0$, defined as a Gaussian random vector of the same size as the pre region:
\begin{equation}
     X^\textrm{pre}_{0} \sim \mathcal{N}(0, I_{N_\textrm{pre}})
\end{equation}
where $N_\textrm{pre}$ is the size of region pre and $\mathcal{N}(0, I_{N_\textrm{pre}})$ is a multivariate Gaussian distribution with 0 mean and identity covariance. This initial random pattern is pattern completed to obtain $X^\textrm{pre}$, and then projected to the post region, to give $X^\textrm{post}$, using in both operations the same fixed semantic load. Then, Hebbian and homeostatic plasticity operations as defined in Algorithms \ref{alg:hebbian} and \ref{alg:homeostasis} are used in the (post, pre) region pair. Replay is computed in simulations following Algorithm \ref{alg:replay}. In the text, we call replay from MTL to CTX \textit{Episodic Replay}, because it finds regularities across recurrent synapses that contain episodic information, and replay from CTX to MTL-semantic \textit{Semantic Replay}, because in this case recurrent synapses contain semantic information.
\begin{algorithm}[h!]
\caption{Replay}\label{alg:replay}
\begin{algorithmic}[1]
\Function{Replay}{post, pre, $S_L$, \texttt{pattern\_complete\_iterations}}
    \State Sample Gaussian random vector in pre:
        \[
            X_0^\textrm{pre} \gets \text{randn}(N_\textrm{pre})
        \]
    \State $X^\textrm{pre} \gets \textsc{PatternComplete}\bigl(X^\textrm{pre}_0, S_L, \texttt{pattern\_complete\_iterations}\bigr)$
    \State $X^\textrm{post} \gets W^{\textrm{post}\leftarrow\textrm{pre}}\cdot X^\textrm{pre} + b^\textrm{post}\textrm{IM}^\textrm{post}$
    \State \textsc{Hebbian}(post, pre)
    \State \textsc{Homeostasis}(post, pre)
    \State \Return $X$
\EndFunction
\end{algorithmic}
\end{algorithm}
\subsection*{Learning Algorithm}
For each day of experience, assuming an input $X_\textrm{input}$ of shape $(N_\textrm{days}, T_\textrm{day}, N_\textrm{SEN})$, we define the following learning algorithm \textsc{Train} (Algorithm \ref{alg:train}), where \textrm{Wake} is defined as in Algorithm \ref{alg:wake} and the \textrm{Sleep} as in Algorithm \ref{alg:sleep}. All operations used have been previously defined, with these algorithms showing the order in which they are performed. \textsc{duration\_phase\_A} represents the number of days in which only episodic replay takes place (Fig. \ref{fig:phase_a}), and  \textsc{duration\_phase\_B} the number of days in which only MTL-sensory synapses are used for episodic replay, but semantic replay starts taking place (Fig. \ref{fig:phase_b}). From day \textsc{duration\_phase\_B} on, MTL replays both sensory and semantic representations (Starting with Fig. \ref{fig:phase_b_supp}). Therefore, during \textsc{duration\_phase\_A} days, episodic memories containing only sensory information are stored in MTL during Wake (Algorithm \ref{alg:wake}). During Sleep (Algorithm \ref{alg:sleep}), cortical abstractions of sensory information are formed. After \textsc{duration\_phase\_A} days, cortical abstractions are transferred to MTL-semantic via semantic replay during Sleep. Finally, after \textsc{duration\_phase\_B} days, episodic memories contain both sensory and semantic information, and so does episodic replay. This allows higher-order abstraction, resulting in the formation of new semantic representations that are based on previous abstractions.
\begin{algorithm}[h!]
\caption{Learning Algorithm Across Days}\label{alg:train}
\begin{algorithmic}[1]
\Function{Train}{$X_\textrm{input}$}
    \For{$\textrm{day} = 1$ to $N_{\textrm{days}}$}
        \State Extract daily input:
            \[
            X^{\textrm{day}}_\textrm{input} \gets X_\textrm{input}[\textrm{day},:]
            \]
        \State \textsc{Wake}($X^{\textrm{day}}_\textrm{input}$)
        \State \textsc{Sleep}()
    \EndFor
\EndFunction
\end{algorithmic}
\end{algorithm}
\begin{algorithm}[h!]
\caption{Wake}\label{alg:wake}
\begin{algorithmic}[1]
\Function{Wake}{$X^{\textrm{day}}$}
        \State Reset MTL recurrent connections:
            \[
                W^{\textrm{MTL} \leftarrow \textrm{MTL}} \gets 0
            \]
    \For{$t = 1$ to $T_{\textrm{day}}$}
        \State $X^{\textrm{SEN}} \gets \textsc{Activation}\bigl(X^{\textrm{day}}[t], \textrm{SEN}\bigr)$
        \State $\hat{X}^{\textrm{MTL-sensory}} 
            \gets 
            W^{\textrm{MTL-sensory} \leftarrow \textrm{SEN}} 
            \cdot X^{\textrm{SEN}}$
        \State $X^{\textrm{MTL-sensory}} 
            \gets 
            \textsc{Activation}\bigl(\hat{X}^{\textrm{MTL-sensory}}, \textrm{MTL-sensory}\bigr)$
        \State $\hat{X}^{\textrm{MTL-semantic}} 
            \gets 
            \text{randn}(N_\textrm{MTL-semantic})$
        \State $X^{\textrm{MTL-semantic}} 
            \gets 
            \textsc{Activation}\bigl(\hat{X}^{\textrm{MTL\_semantic}}, \textrm{MTL\_semantic}\bigr)$
        \State $X^{\textrm{MTL}} 
            \gets 
            [X^{\textrm{MTL-sensory}}, \; X^{\textrm{MTL-semantic}}]$
        \State $\hat{X}^{\textrm{CTX}} 
            \gets 
            W^{\textrm{CTX} \leftarrow \textrm{MTL}} 
            \cdot X^{\textrm{MTL}}$
        \State $X^{\textrm{CTX}} 
            \gets 
            \textsc{Activation}\bigl(\hat{X}^{\textrm{CTX}}, \textrm{CTX}\bigr)$
        \If{day $\ge$ duration\_phase\_A}
            \State $\hat{X}^{\textrm{MTL-semantic}} 
                \gets 
                W^{\textrm{MTL-semantic} \leftarrow \textrm{CTX}} 
                \cdot X^{\textrm{CTX}}$
            \State $X^{\textrm{MTL-semantic}} 
                \gets 
                \textsc{Activation}\bigl(\hat{X}^{\textrm{MTL-semantic}}, \textrm{MTL-semantic}\bigr)$
            \State $X^{\textrm{MTL}}[N_\textrm{MTL-sensory}:] 
                \gets 
                X^{\textrm{MTL\_semantic}}$
            \State \textsc{Hebbian}$(X^{\textrm{MTL\_semantic}}, X^{\textrm{MTL\_semantic}})$
        \EndIf
        \If{day $\ge$ duration\_phase\_B}
            \State $\hat{X}^{\textrm{CTX}} 
                \gets 
                W^{\textrm{CTX} \leftarrow \textrm{MTL}} 
                \cdot X^{\textrm{MTL}}$
            \State $X^{\textrm{CTX}} 
                \gets 
                \textsc{Activation}\bigl(\hat{X}^{\textrm{CTX}}, \textrm{CTX}\bigr)$
        \EndIf
        \State \textsc{Hebbian}$(\textrm{MTL}, \textrm{MTL})$
        \State \textsc{Hebbian}$(\textrm{CTX}, \textrm{CTX})$
        \State \textsc{Homeostasis}$(\textrm{CTX}, \textrm{CTX})$
    \EndFor
    \State Increment day counter
\EndFunction
\end{algorithmic}
\end{algorithm}
\begin{algorithm}[h!]
\caption{Sleep}\label{alg:sleep}
\begin{algorithmic}[1]
\Function{Sleep}{}
\State Episodic Replay
    \For{$t = 1$ to $\texttt{sleep\_duration\_A}$}
        \If{day $\le$ duration\_phase\_B}
            \State \textsc{Replay}(\textrm{CTX}$, \textrm{MTL-sensory}$, $S_L= 1$, $\textrm{pattern\_complete\_iterations}=10$)
        \Else
        \State Sample $S_L$ randomly
        \State $S_L \gets \textsc{randint}(1, S_{L-\textrm{max}})$
            \State \textsc{Replay}(\textrm{CTX}$, \textrm{MTL}$, $S_L= S_L$, $\textrm{pattern\_complete\_iterations}=10$)
        \EndIf
    \EndFor

    \If{day $\ge$ duration\_phase\_A}
        \For{$t = 1$ to $\texttt{sleep\_duration\_B}$}
            \State \textsc{Replay}(\textrm{MTL-semantic}$, \textrm{CTX}$, $S_L= 1$, $\textrm{pattern\_complete\_iterations}=10$)
        \EndFor
    \EndIf
\EndFunction
\end{algorithmic}
\end{algorithm}
\subsection*{Considerations regarding modeling choices}
\begin{itemize}
    \item Input generation
    \item top-K
    \item Outgoing vs incoming homeostasis
    \item Neuronal maturity
    \item Sleep time
\end{itemize}
\clearpage
\subsection*{Data Analysis}
\subsubsection*{Neuronal Selectivity}
Given a vector $X$ that contains neuronal activity across T timesteps -with dimensions $(T, N_\textrm{neurons})$- and one that contains certain latent variables $Z$ -with dimensions  $(T, N_\textrm{latents})$, we do obtain the normalized values:
\begin{equation}
    X_\textrm{norm} = \frac{X - \langle X \rangle_{\textrm{time}}}{\sqrt{\langle X^2 \rangle_\textrm{time} - \langle X \rangle_\textrm{time}^2}} \;\; ; \;\; Z_\textrm{norm} = \frac{Z - \langle Z \rangle_{\textrm{time}}}{\sqrt{\langle Z^2 \rangle_\textrm{time} - \langle Z \rangle_\textrm{time}^2}}
\end{equation} 
And then obtain the matrix 
\begin{equation}
    S = \frac{1}{T}X_\textrm{norm}^TZ_\textrm{norm}
\end{equation}
which, for neuron $i$ and latent $j$ contains their correlation as:
\begin{equation}
    S(i, j) = \frac{1}{T}\sum_t X_{ti}Z_{tj}
\end{equation}
The Maximum Selectivity of neuron $i$ is the maximum selectivity value of neuron $i$ across all latents $j$:
\begin{equation}
    S_\textrm{max}(i) = \max_{j}{S(i, j)}
\end{equation}
\subsubsection*{Neuron Ordering}
Throughout this study, we re-order the neurons indices for visualization purposes. We do this with respect to a certain latent variable (for example the different \textit{concepts} $A_i$ and $B_i$ or the different episodes $A_iB_j$. This is done by imposing a threshold, and then for each possible discrete value the latent can take we obtain the neurons with a selectivity higher to that value with respect to the given latent. For example, assuming latents $A_1$ to $B_5$, the first neurons would be those with a selectivity higher than 0.9 to $A_1$, then the next neurons would be those selective to $A_2$, etc. All neurons that are not selective to any latent keep the original order at the end of the new order.
\subsubsection*{Signal-to-Noise Ratio}
In order to obtain the Signal-to-Noise Ratio (SNR) of a generated pattern, we first define the signal  $\mathcal{S}$  as the expected overlap (number of common neurons) of the corrupted pattern with the original, and subtract the expected overlap between two randomly generated patterns of $K^\textrm{region}$ active neurons across $N^\textrm{region}$ neurons:
\begin{equation}
\mathcal{S}  =  \textrm{E}\Big[|| X^\textrm{corrupted} \odot X^\textrm{original}||_0\Big] - \textrm{E}\Big[|| X^\textrm{random-1} \odot X^\textrm{random-2}||_0\Big]
\end{equation}
Because we fix that $N_\textrm{swaps}$ active (in the original pattern) flip from 1 to 0, the first term is always $K^\textrm{region} - N_\textrm{swaps}$, while the second corresponds to the expected value of a hypergeometric distribution, given by $(K^\textrm{region})^2/N^\textrm{region}$. Then, we define the noise to be the number of non-common neurons between the corrupted and the original pattern, which is $2N_\textrm{swaps}$, giving:
\begin{equation}
    \mathrm{SNR} = \frac{\Big(K - N^\textrm{swaps}- K^2/N\Big)^2}{\Big(2N^\textrm{swaps}\Big)^2}
    \label{eq:snr}
\end{equation}
\subsection*{Probability of adding concepts to cortex}
To quantify how exposure statistics shape cortical dictionary formation (Figs. \ref{fig:phase_b}h and \ref{fig:higher-order}f), we measured, across many independently simulated networks, whether a targeted simple concept or a targeted full episode became to the cortical dictionary after a session of sleep as a function of its probability of occurrence in the training distribution. In both analyses, each simulation produced a binary formation outcome for one focal item, and the plotted quantity was the empirical probability of formation obtained by averaging this binary outcome across all simulations sharing the same focal probability.

\paragraph*{Simple concepts (Fig. \ref{fig:phase_a}h).}
For the simple-concept analysis, we simulated one day of awake input ($T_{\mathrm{day}}=80$, mean episode duration $5$, fixed duration, input corruption $N_{\mathrm{swaps}}=4$) followed by Sleep A only ($10$ sleep timesteps, no Sleep B) in newly initialized networks. In each simulation, one concept was selected independently within each latent dimension, and its marginal probability was set by sampling from a discrete grid
\[
\{0.1,0.2,\dots,0.9\}.
\]
The remaining probability mass within that latent dimension was distributed across the other concepts by Dirichlet sampling, and episode probabilities were then constructed as the product of the two latent marginals,
\begin{equation}
p(A_i,B_j)=p(A_i)\,p(B_j).
\end{equation}
After training and sleep, we identified mature CTX neurons as those with $\mathrm{ctx\_IM}=0$, and assigned each such neuron to the simple concept whose MTL-sensory input weights had the largest mean value over the neurons belonging to that concept. A concept was counted as \emph{added to the cortical dictionary} if at least one mature CTX neuron was assigned to it. For each simulation, we recorded the pair
\begin{equation}
\big(p_{\mathrm{focal}},\, \mathbbm{1}_{\mathrm{formed}}\big),
\end{equation}
where $p_{\mathrm{focal}}$ is the marginal probability of the focal concept and $\mathbbm{1}_{\mathrm{formed}}\in\{0,1\}$ indicates whether it was represented. We repeated this analysis across $1000$ seeds, grouped identical focal probabilities, and plotted the mean of $\mathbbm{1}_{\mathrm{formed}}$ at each probability level.

\paragraph*{Complex episodes (Fig. \ref{fig:higher-order}f).}
For the complex concepts, we started from the pretrained phase-B checkpoint and evaluated one additional day of training with $S_L^\textrm{max}=2$, the same awake input structure ($T_{\mathrm{day}}=80$, mean episode duration $5$, input corruption $N_{\mathrm{swaps}}=8$), followed by sleep in the intact network. In each simulation, one focal episode $(A_i,B_j)$ was chosen uniformly from the $25$ possible episodes, and its probability of occurrence in training was sampled from a discrete grid
\[
\{0,0.025,0.05,\dots,0.5,0.6,0.7,0.8,0.9,1.0\}.
\]
The focal episode was assigned this probability, and the remaining mass was distributed uniformly across the other $24$ episodes,
\begin{equation}
p(A_k,B_\ell)=
\begin{cases}
p_{\mathrm{focal}}, & (k,\ell)=(i,j),\\[4pt]
\dfrac{1-p_{\mathrm{focal}}}{24}, & \text{otherwise.}
\end{cases}
\end{equation}
After training, we evaluated whether the focal episode had been incorporated into the higher-order cortical dictionary. To do so, we extracted the receptive fields from the episode-selective CTX subregion and compared them with the ideal episode prototypes in MTL-semantic space using cosine similarity. For each episode prototype $e$, we computed
\begin{equation}
s_e=\max_n \cos\!\big(\mathbf{p}_e,\mathbf{w}_n\big),
\end{equation}
where $\mathbf{p}_e$ is the prototype of episode $e$ and $\mathbf{w}_n$ is the receptive field of CTX neuron $n$. The focal episode was considered \emph{added to CTX} when its maximal similarity exceeded a threshold of $0.8$. As above, each simulation yielded a pair
\begin{equation}
\big(p_{\mathrm{focal}},\, \mathbbm{1}_{\mathrm{formed}}\big),
\end{equation}
and the plotted curve was obtained by averaging $\mathbbm{1}_{\mathrm{formed}}$ across all simulations with the same focal episode probability. We repeated this analysis across $2000$ seeds.
\subsection*{Impact of semantic representations in episodic memory}
We compared an intact network that contained semantic representations in MTL \emph{semantics present} with two lesioned networks: one without MTL-semantic (\emph{semantics absent}) and one with MTL-semantic neural activity randomized during encoding (to control for the total number of neurons available, (\emph{semantics random}).
\subsection*{Compositional generalization and replay}
To quantify how semantic representations in episodic memory shape cortical compositional learning, we examine learning of \textit{complex concepts}. This is allowed by increasing the maximum \textit{semantic load} during replay to 2, and performed across wide range of input noise levels. Input noise levels are again controlled by $N_\textrm{swaps}$, with each $N_\textrm{swaps}$ mapping to a particular Signal-to-Noise Ratio (SNR) via Equation \eqref{eq:snr}. 
For each noise level and seed, networks were trained with the rest of parameters intact as in Figs. \ref{fig:phase_a}-\ref{fig:phase_b}, with the last $N^{S_L=2, \mathrm{train}}_\textrm{days}$ training days involving a semantic load of 2, and used to determine cortical neuronal selectivity to each complex concept $A_iB_j$. Then, cortical activity was used to decode episodic identity across $N^{S_L=2, \mathrm{test}}_\textrm{days}$ days, from which the Compositional Generalization (Accuracy) measure was obtained.
\newline\newline
To estimate the effective SNR of replayed patterns (Fig. \ref{fig:normative}c), we analyze Sleep A replays with semantic load of 2. Each replayed MTL-sensory pattern is matched to its closest de-noised version (MTL-sensory patterns generated with $N_\textrm{swaps}=0$), and the Hamming distance to thhis noiseless pattern is obtained (number of neurons with different activity). Then, the $N_\textrm{swaps}$ that would result in such a number of neurons with different activity is calculated dividing by 2. This allows us to have an estimate of the SNR associated with replayed MTL-sensory pattern by inverting which SNR of input generation could result in a pattern so distant from noiseless patterns. It should be noted, however, these patterns are completely internally generated from random activity.
\newline\newline
To quantify Replayed / Awake Generalization (Fig. \ref{fig:normative}d), we compare Sleep A replayed MTL-sensory patterns to a separate set of noisy awake input patterns generated under the same corruption level. Each replay is first assigned a semantic label by matching its MTL-sensory state to the closest noiseless pattern. For each awake sample, we then compute its mean cosine similarity to all replayed patterns whose inferred label $A_iB_j$ corresponds to the awake sample, and  subtract from it the mean cosine similarity with the rest of replayed patterns. The average of this difference across awake patterns is defined as the Replayed / Awake Generalization. Only MTL-sensory is used because during awake MTL-semantic is first inferred from MTL-sensory (and therefore is not involved in CTX selectivity - unless additional cortical noise is injected as in Fig. \ref{fig:behavioural}g). The Replayed / Awake generalization score can be interpreted as \textit{For each replayed pattern, if it was used as a receptive field, would it correctly identify whether awake input patterns represent the same episode?}.
\newline\newline
To quantify Replayed / Replayed generalization (Fig. \ref{fig:normative}e), we compare full MTL replay states to one another during Sleep A. Each replay is again assigned a semantic label from its MTL-sensory component by matching it to the closest noiseless pattern. We then compute the cosine similarity matrix between all replayed (full) MTL patterns. For each replay event, we compare its mean similarity to other replays with the same inferred label against its mean similarity to replays with different labels. In this case, the CTX activity that will drive plasticity is affected by all of MTL, which is why the cosine similarity is obtained from the full MTL representation. The Replayed / Replayed generalization score can be interpreted as \textit{For each replayed pattern, if it was used as a receptive field, would it correctly identify whether the rest of replayed patterns represent the correct episode?}.
\subsection*{Recall capacity}
To quantify recall capacity (Fig. \ref{fig:higher_order}a-b), we used a common recall assay across all conditions. A set of patterns was first stored in the recurrent memory system under a given input corruption level, defining the signal-to-noise ratio of the stored representations SNR (stored). Each stored pattern was then probed with a noisy cue obtained by applying a fixed number of swaps, corresponding to cue corruption SNR (cue). Starting from this corrupted cue, recurrent pattern completion was run for a fixed number of iterations in the relevant memory system. Recall was quantified as the cosine similarity between the completed pattern and the original stored target, and was then averaged across stored patterns and across seeds.
\newline\newline
As a baseline benchmark (Fig. \ref{fig:higher_order}a), we first measured recall capacity for random patterns. We compared three associative-memory systems: the MTL-sensory subregion alone, the full MTL recurrent system, and a sparse Hopfield network with matched sparsity constraints. In each case, increasing numbers of random sparse patterns were stored directly in the corresponding recurrent weights, and recall was then tested with the common noisy-cue assay described above. This analysis provides a classical capacity benchmark that is independent of semantic or compositional structure.
\newline\newline
We next turned to structured episodic patterns generated from the standard compositional latent space (Fig. \ref{fig:higher_order}b). Instead of storing random binary patterns, networks were exposed to inputs generated from pairs $(A_i,B_j)$ drawn from the within-distribution environment, and recall was measured as the number of stored patterns increased. We compared the same semantic conditions used throughout the manuscript: \emph{semantics present}, in which recall uses the intact episodic-semantic circuit; \emph{semantics absent}, in which recall is restricted to sensory recurrence; and \emph{semantics random}, in which semantic units remain present but are set to random for each episode. Unless otherwise noted, recall was quantified in MTL-sensory; when indicated explicitly, we also report recall within MTL-semantic (taking the cosine similarity across MTL-subpopulation).
\newline\newline
To interpret these recall-capacity differences, we compared the effective pattern entropy of the different stored representations. For random sparse patterns with $N$ neurons and exactly $K$ active units, the number of possible patterns is
\begin{equation}
\Omega_{\mathrm{random}} = \binom{N}{K},
\end{equation}
yielding entropy
\begin{equation}
H_{\mathrm{random}} = \log_2 \binom{N}{K}.
\end{equation}
For structured MTL-semantic patterns, entropy is approximated by the entropy of episode identity,
\begin{equation}
H_{\mathrm{MTL\text{-}semantic}} = H(A,B)
= - \sum_{i,j} p_{ij}\log_2 p_{ij},
\end{equation}
where $p_{ij}=p(A=i,B=j)$ denotes the episode distribution. For structured MTL-sensory patterns, entropy includes both episode identity and the additional combinatorial degrees of freedom introduced by sensory corruption,
\begin{equation}
H_{\mathrm{MTL\text{-}sensory}}
=
H(A,B)
+
\log_2 \binom{K}{N_{\mathrm{swaps}}}
+
\log_2 \binom{N-K}{N_{\mathrm{swaps}}},
\end{equation}
where $N_{\mathrm{swaps}}$ is the number of active units turned off and inactive units turned on. This theoretical comparison provides a compact way of relating recall capacity to the statistical complexity of the stored patterns, rather than to circuit architecture alone.
\subsection*{Simulation parameters}
Default input, temporal, and model parameters can be found in Tables \ref{tab:default_input}-\ref{tab:default_model}. Figure-specific parameters can be found in Tables \ref{tab:figure_2}-\ref{tab:figure_6}.
\begin{table}[ht]
\centering
\caption{Summary of Episode Generation Protocol parameters}
\label{tab:default_input}
\begin{tabular}{lll}
\toprule
\textbf{Variable} & \textbf{Description} & \textbf{Value} \\
\midrule
$N_\textrm{attributes}$ & Number of episode attributes & 2 \\
$|A|, |B| $ & Number of concepts per attribute & [5, 5] \\
$N_\textrm{concept}^A$, $N_\textrm{concept}^B$ & Total size (in neurons) per concept & [10, 10] \\
$P(A_i, B_j)$ & Episode probability & $0.5/5$ if $i = j$ ; $0.5/20$ if $i\neq j$ \\
$N_\textrm{swaps}$ & Number of swaps (total noise is $2N_\textrm{swaps}$) & 4 \\
\bottomrule
\end{tabular}
\end{table}
\begin{table}[ht]
\centering
\caption{Summary of default input temporal parameters}
\label{tab:default_time}
\begin{tabular}{lll}
\toprule
\textbf{Variable} & \textbf{Description} & \textbf{Value} \\
\midrule
$T_\textrm{day}$ & Day length & 80 \\
$T_\textrm{episode}$ & Episode length & 5 \\
\bottomrule
\end{tabular}
\end{table}

\begin{table}[ht]
\small
\centering
\caption{Summary of network parameters used in the simulations (default unless specified otherwise).}
\label{tab:default_model}
\begin{tabular}{lp{6cm}l}
\toprule
\textbf{Variable} & \textbf{Description} & \textbf{Value} \\
\midrule
\multicolumn{3}{l}{\textbf{General}} \\
\midrule
$\texttt{duration\_phase\_A}$ & Phase A duration (episodic $\to$ semantic only) & 1000 \\
$\texttt{duration\_phase\_B}$ & Phase B duration (episodic $\to$ semantic + semantic $\to$ episodic) & 1500 \\
$\texttt{sleep\_duration\_episodic}$ & Sleep duration for episodic replay & 10 \\
$\texttt{sleep\_duration\_semantic}$ & Sleep duration for semantic replay & 10 \\
$S_L^\textrm{min}$ & Min semantic load during replay & 1 \\
$S_L^\textrm{max}$ & Max semantic load during input & 2 \\
\midrule
\multicolumn{3}{l}{\textbf{Pattern Completion}} \\
\midrule
$\texttt{mtl\_pattern\_complete\_iterations}$ & Iterations in MTL pattern completion & 10 \\
$\texttt{ctx\_pattern\_complete\_iterations}$ & Iterations in CTX pattern completion & 10 \\
\midrule
\multicolumn{3}{l}{\textbf{Region Parameters}} \\
\midrule
$N_{\text{sen}}$ & Number of SEN neurons & 100 \\
$K^\textrm{SEN}_\textrm{concept}$ & Neurons per concept in SEN & 10 \\
$K^\textrm{SEN}_\textrm{episode}$ & Neurons per episode in SEN & 20 \\
$N^{\text{CTX}}$ & Number of CTX neurons (list denotes sub-regions) & [100, 300] \\
$K^\textrm{CTX}_\textrm{concept}$ & Neurons per concept in CTX (list denotes sub-regions) & [10,  10] \\
$K^\textrm{CTX}_\textrm{episode}$ & Neurons per episode in CTX (list denotes sub-regions) & [20,  10] \\
$b^\text{CTX}$ & CTX immature neurons excitability (list denotes sub-regions) & [0.3, 0.6] \\
$N_{\text{MTL-sensory}}$ & Number of MTL-sensory neurons & 100 \\
$K^\textrm{MTL-sensory}_\textrm{concept}$ & Neurons per concept in MTL-sensory & 10 \\
$K^\textrm{MTL-sensory}_\textrm{episode}$ & Neurons per episode in MTL-sensory & 20 \\
$N_{\text{MTL-semantic}}$ & Number of MTL-semantic neurons & 100 \\
$K^\textrm{MTL-semantic}_\textrm{concept}$ & Neurons per concept in MTL-semantic & 5 \\
$K^\textrm{MTL-semantic}_\textrm{episode}$ & Neurons per episode in MTL-semantic & 10 \\
$b^\text{MTL.semantic}$ & MTL-semantic immature neurons excitability & 0.3 \\
\midrule
\multicolumn{3}{l}{\textbf{Synaptic Parameters}} \\
\midrule
$W_\sigma^{\text{CTX}\leftarrow\text{MTL}}$ & CTX$\leftarrow$MTL weight init std &$5\cdot 10^{-4}$ \\
$\lambda^{\text{CTX}\leftarrow\text{MTL}}_{\textrm{slow}}$ & CTX$\leftarrow$MTL learning rate & $5\cdot 10^{-4}$ \\
$W^{\text{CTX}\leftarrow\text{MTL}}_{\text{max-pre}}$ & Max outgoing CTX$\leftarrow$MTL weight & $\infty$ \\
$W^{\text{CTX}\leftarrow\text{MTL}}_{\text{max-post}}$ & Max incoming CTX$\leftarrow$MTL weight & 1 \\
$\lambda^{\text{MTL}\leftarrow\text{MTL}}_0$ & MTL recurrent learning rate & $5\cdot 10^{-3}$ \\
$W^{\text{MTL}\leftarrow\text{MTL}}_{\text{max-pre}}$ & Max outgoing MTL$\leftarrow$MTL weight & $\infty$ \\
$W^{\text{MTL}\leftarrow\text{MTL}}_{\text{max-post}}$ & Max incoming MTL$\leftarrow$MTL weight & $\infty$ \\
$\lambda^{\text{CTX}\leftarrow\text{CTX}}_0$ & CTX recurrent learning rate & $5\cdot 10^{-4}$ \\
$W^{\text{CTX}\leftarrow\text{CTX}}_{\text{max-pre}}$ & Max outgoing CTX$\leftarrow$CTX weight & 1 \\
$W^{\text{CTX}\leftarrow\text{CTX}}_{\text{max-post}}$ & Max incoming CTX$\leftarrow$CTX weight & $\infty$ \\
$W_\sigma^{\text{MTL-semantic}\leftarrow\text{CTX}}$ & MTL-semantic$\leftarrow$CTX weight init std & $1\cdot 10^{-3}$ \\
$\lambda^{\text{MTL-semantic}\leftarrow\text{CTX}}_{\textrm{slow}}$ & MTL-semantic$\leftarrow$CTX learning rate & $5\cdot 10^{-3}$ \\
$W^{\text{MTL-semantic}\leftarrow\text{CTX}}_{\text{max-pre}}$ & Max outgoing MTL-semantic$\leftarrow$CTX weight & $\infty$ \\
$W^{\text{MTL-semantic}\leftarrow\text{CTX}}_{\text{max-post}}$ & Max incoming MTL-semantic$\leftarrow$CTX weight & 1 \\
\bottomrule
\end{tabular}
\end{table}
\begin{table}[ht]
\centering
\caption{Figure 2 parameters (different than or not specified in default)}
\label{tab:figure_2}
\begin{tabular}{lll}
\toprule
\textbf{Variable} & \textbf{Description} & \textbf{Value} \\
\midrule
$\texttt{MTL-sensory\_SEN\_projection}$ & MTL-sensory$\leftarrow$SEN random projection & False \\
$N_\textrm{days}$ & Number of days & 1000 \\
$N_\textrm{swaps}$ & Number of swaps (total noise is $2N_\textrm{swaps}$) & 4 \\
\bottomrule
\end{tabular}
\end{table}

\begin{table}[ht]
\centering
\caption{Figure 3 parameters (different than or not specified in default)}
\label{tab:figure_3}
\begin{tabular}{lll}
\toprule
\textbf{Variable} & \textbf{Description} & \textbf{Value} \\
\midrule
\makecell{$\texttt{MTL-sensory}$\\$\texttt{\_SEN\_projection}$} & MTL-sensory$\leftarrow$SEN random projection & False \\
$N_\textrm{swaps}$ & Number of swaps (total noise is $2N_\textrm{swaps}$) & 4 \\
\midrule
Figs. 3b-3e & & \\
$N_\textrm{days}$ & Number of days & 1500 \\
\midrule
Fig. 3g-3i & & \\
$N_\textrm{Days}$ & Number of days (to get classifier train sample, after model trained) & 100 \\
\midrule
Fig. 3h & & \\
$P(A_i, B_j)$ & Episode probability (train) & 0 if $i=j$ ; 1/20 if $i\neq j$ \\
$P(A_i, B_j)$ & Episode probability (test) & 1/5 if $i=j$ ; 0 else \\
\bottomrule
\end{tabular}
\end{table}


\begin{table}[ht]
\centering
\caption{Figure 4 parameters (different than or not specified in default)}
\label{tab:figure_4}
\begin{tabular}{lll}
\toprule
\textbf{Variable} & \textbf{Description} & \textbf{Value} \\
\midrule
Interleaved (pre-training) \\
$N_\textrm{days}$ & Number of days & 100 \\
$T_\textrm{day}$ & Day length & 200 \\
$T_\textrm{episode}$ & Episode length & 1 \\
\midrule
Blocked (pre-training) \\
$N_\textrm{days}$ & Number of days & 10 (per concept)\\
$T_\textrm{day}$ & Day length & 200  \\
$T_\textrm{episode}$ & Episode length & 1 \\
$P(A_i, B_j)$ & Episode probability & 0.2 if block concept present (0 else) \\
\midrule
Training \\
$N_\textrm{days}$ & Number of days & 1400 \\
\bottomrule
\end{tabular}
\end{table}

\begin{table}[ht]
\centering
\caption{Figure 5 parameters (different than or not specified in default)}
\label{tab:figure_5}
\begin{tabular}{lll}
\toprule
\textbf{Variable} & \textbf{Description} & \textbf{Value} \\
\midrule
Fig. 5c and 5e (Start with best accuracy network from Fig. 4)\\
$N_\textrm{days}$ & Number of days & 1 \\
$T_\textrm{day}$ & Day length & 25 \\
$T_\textrm{episode}$ & Episode length & 5 \\
\midrule
Fig. 5d (Start with best accuracy network from Fig. 4)\\
$N_\textrm{days}$ & Number of days & 1 \\
$T_\textrm{day}$ & Day length & 1 to 5000 (spaced) \\
$T_\textrm{episode}$ & Episode length & 5 \\
\midrule
Fig. 5f \\
\midrule 
Sensory replayed \\
1. Pre-train with blocked (see Table \ref{tab:figure_4})\\
2. 1400 days\\
3. 100 days frozen to test accuracy \\
\midrule
Sensory-semantic replayed \\
$S_L^\textrm{max}$ & \makecell{Max. semantic \\ load during replay} & 1 \\
1. Pre-train with blocked (see Table \ref{tab:figure_4})\\
2. 2000 days \\
3. 100 days frozen to test accuracy \\

\bottomrule
\end{tabular}
\end{table}

\begin{table}[ht]
\centering
\caption{Figure 6 parameters (different than or not specified in default)}
\label{tab:figure_6}
\begin{tabular}{lll}
\toprule
\textbf{Variable} & \textbf{Description} & \textbf{Value} \\
\midrule
Start with best accuracy network from Fig. 5f \\
$N_\textrm{days}$ & Number of days & 500 \\
\bottomrule
\end{tabular}
\end{table}

\subsection*{Reproduction of experimental results}
\subsubsection*{Concept cell formation \shortcite{quiroga_2005, quiroga_2012}}
\paragraph{Experimental goal.}
We quantified concept-cell formation via different measures aimed at quantifying their main defining properties: a high selectivity to individual concepts that generalizes across contexts and metrics. To do so, we measure (i) cross-context linear decodability, (ii) concept selectivity and (iii) normalized mutual information.

\paragraph{Network and latent space.}
We used a pre-trained network (same training regime as Figure 3) and evaluated activity in MTL-sensory and MTL-semantic. As before, episodes were defined as the conjunction of two concepts $(A_i, B_j)$.

\paragraph{Cross-context linear decodability of concept representations}
To quantify whether individual concepts are represented in a way that generalizes across episodic contexts (Fig. \ref{fig:circuits}a), we used a cross-context linear decoding assay. We used a pre-trained network (same training regime as Fig. \ref{fig:phase_b}) and evaluated activity in MTL-sensory and MTL-semantic under frozen weights. As before, episodes were defined as conjunctions of two concepts $(A_i,B_j)$, with $A_i \in \{A_1,\dots,A_5\}$ and $B_j \in \{B_1,\dots,B_5\}$.
\newline\newline
For each target concept $c \in A \cup B$, we constructed a binary classification task. If the target was $A_i$ ($c = A_i$), we first sampled one concept $B_j$ to define a \textit{training context}; if the target was $B_j$, we analogously sampled one concept $A_i$. Training data then consisted of exactly two episode types presented in one deterministic day: one positive episode containing the target concept, and one negative episode in which the target concept was replaced by a contrast concept from the same attribute. Thus, for a target $A_i$ and training context $B_j$, the two training episodes were $(A_i,B_j)$ and $(A_k,B_j)$ with $k \neq i$. Likewise, for a target $B_j$ and training context $A_i$, the two training episodes were $(A_i,B_j)$ and $(A_i,B_\ell)$ with $\ell \neq j$. Each episode was presented as one contiguous block of length $T_\textrm{episode}$, yielding a minimal balanced training set by construction.
\newline\newline
From this deterministic training day, we extracted awake activity from MTL-sensory and MTL-semantic, denoted $X^\textrm{MTL-sensory}$ and $X^\textrm{MTL-semantic}$. Labels were defined as
\[
Y_c =
\begin{cases}
1 & \text{if concept } c \text{ was present in the episode} \\
0 & \text{otherwise.}
\end{cases}
\]
For each region separately, we then trained a linear logistic classifier to predict $Y_c$ from the corresponding neural activity. In all cases, the same labels were used for both regions, so differences in decoding performance reflected only differences in the underlying neural representations.
\newline\newline
To test cross-context generalization, we generated a second dataset in which the training context was excluded. Concretely, if the decoder had been trained to recognize $A_i$ in the context of $B_j$, then all test episodes containing $B_j$ were removed from the sampling distribution. The remaining probability mass was redistributed so that the marginal probability of the target concept remained 0.5, and the remaining 0.5 was distributed uniformly across non-target concepts. The same logic was applied symmetrically when decoding a target $B_j$. This yielded a test distribution in which the target concept could only be identified if its representation generalized beyond the specific opposite-attribute context used during training.
\newline\newline
For each target concept, we repeated this procedure across multiple random choices of training context and contrast concept. Decoder performance was quantified as held-out classification accuracy on the excluded-context test set. The final cross-context linear decodability score for each region was obtained by aggregating these accuracies across repeats and across all target concepts in $A \cup B$. In Fig. \ref{fig:circuits}a, we compare the resulting concept-level distributions for MTL-sensory and MTL-semantic.
\paragrpah*{Latent selectivity of MTL neurons}
To quantify how selectively individual MTL neurons represented latent concepts (Fig. \ref{fig:circuits}b), we measured the selectivity of neurons in MTL-sensory and MTL-semantic to the simple concepts $A_i$ and $B_j$. To do so, we used the same pre-trained network as in Fig. \ref{fig:phase_b}, froze all synaptic weights, and generated a long test dataset with uniform sampling over all possible episodes $(A_i,B_j)$. This ensured that selectivity was estimated under a controlled and unbiased distribution over latent combinations.
\newline\newline
Specifically, we sampled one long test day of length $T_\textrm{day}=10000$ timesteps, with fixed episode duration and the same level of within-episode swap noise used elsewhere in the manuscript. Episode probabilities were set to
\[
P(A_i,B_j)=\frac{1}{|A||B|},
\]
so that all $25$ possible pairs were equally likely. The frozen network was then run on this input while recording awake activity in MTL-sensory and MTL-semantic. This yielded activity matrices
\[
X^\textrm{MTL-sensory}\in\mathbb{R}^{T\times N_\textrm{MTL-sensory}},
\qquad
X^\textrm{MTL-semantic}\in\mathbb{R}^{T\times N_\textrm{MTL-semantic}},
\]
where $T$ is the number of awake samples.
\newline\newline
To define the latent variables, we converted the episode labels into one-hot concept indicators for the two attributes separately. Thus, for each timestep we constructed a binary latent vector
\[
Z_t = (Z_t^{A_1},\dots,Z_t^{A_5},Z_t^{B_1},\dots,Z_t^{B_5}),
\]
with one active entry in the $A$ block and one active entry in the $B$ block. We then quantified the selectivity of neuron $n$ to latent concept $\ell$ as the normalized correlation between its activity trace and the corresponding latent indicator:
\[
s_{n\ell}
=
\frac{1}{T}
\sum_{t=1}^{T}
\tilde{X}_{tn}\tilde{Z}_{t\ell},
\]
where $\tilde{X}_{tn}$ and $\tilde{Z}_{t\ell}$ denote zero-mean, variance-normalized versions of neural activity and latent variables, respectively. In this way, $s_{n\ell}$ measures how strongly neuron $n$ is associated with the presence of concept $\ell$ across the test dataset.
\newline\newline
For each neuron, we retained its maximal latent selectivity
\[
s_n^\ast = \max_{\ell} s_{n\ell},
\]
which provides a single measure of how selectively that neuron represents any simple concept. We then restricted the analysis to highly selective neurons by applying a threshold to $s_n^\ast$, and compared the resulting distributions between MTL-sensory and MTL-semantic. In Fig. \ref{fig:circuits}b, this yielded a region-wise comparison of the strength of concept selectivity at the single-neuron level.

\paragraph*{Mutual information of selective neurons}
To complement the selectivity analysis (Fig. \ref{fig:circuits}b), we next quantified how much information highly selective neurons carried about their preferred latent concept (Fig. \ref{fig:circuits}c). This analysis used the same long frozen test rollout and the same latent variables as above, so that selectivity and mutual information were estimated from identical neural activity under the same uniform episode distribution.
\newline\newline
Starting from the selectivity matrix $\{s_{n\ell}\}$, we first identified neurons whose maximal latent selectivity exceeded a fixed threshold. For each such neuron $n$, we then assigned a \textit{preferred latent concept}
\[
\ell_n^\ast = \arg\max_{\ell} s_{n\ell},
\]
that is, the concept for which the neuron showed the strongest selectivity. The goal of the mutual-information analysis was then to quantify how informative the activity of neuron $n$ was about the binary variable indicating the presence or absence of its preferred concept $\ell_n^\ast$.
\newline\newline
For each neuron, we therefore considered the binary latent variable
\[
Z^{(\ell_n^\ast)}_t \in \{0,1\},
\]
which indicates whether the preferred concept of that neuron was present at timestep $t$. We then computed the mutual information between the neuron activity trace $X_n$ and this binary latent variable,
\[
I(X_n; Z^{(\ell_n^\ast)}),
\]
using the empirical joint distribution across awake samples in the long test day. To facilitate comparison across neurons, this quantity was normalized by the entropy of the preferred latent variable,
\[
\mathrm{NMI}_n
=
\frac{I(X_n; Z^{(\ell_n^\ast)})}{H(Z^{(\ell_n^\ast)})}.
\]
This normalization yields a value between 0 and 1, where 0 indicates that the neuron carries no information about the presence of its preferred concept, and 1 indicates that its activity fully determines that concept variable.
\begin{table}[ht]
\centering
\caption{Parameters used for concept-cell formation and selectivity analyses.}
\label{tab:concept_cell_parameters}
\begin{tabular}{p{0.58\linewidth}p{0.30\linewidth}}
\toprule
\textbf{Parameter} & \textbf{Value} \\
\midrule
Latent attributes and concepts & 2 attributes ($A$, $B$), 5 concepts each \\
Concept code size (object, scene) & 10 active units, 10 active units \\
Training setup per concept task & 1 fixed opposite context, 1 target concept, 1 contrast concept \\
Training class structure & Two balanced blocks (positive vs negative) \\
Training block duration & 1 timestep per block \\
Training swap noise & 8 swap operations per timestep \\
Test sampling for decoding & 1 day, 200 timesteps, fixed episode duration 5 \\
Test context rule & Exclude all episodes containing training opposite-context concept \\
Target prevalence in test set & 0.5 (remaining mass uniform over non-target cases) \\
Concept repeats & 10 repeats per concept \\
Decoder type & Linear logistic regression (L2 regularization) \\
Regions decoded & MTL-sensory, MTL-semantic \\
Long selectivity test stream & 1 day, 10,000 timesteps, uniform over all latent pairs \\
Selectivity threshold (reported notebook run) & 0.6 \\
Mutual-information readout & $I(x_n;z_{\ell^\*})/H(z_{\ell^\*})$ on thresholded units \\
\bottomrule
\end{tabular}
\end{table}
\subsubsection*{Consolidation promotes the emergence of representational overlap in the hippocampus and medial prefrontal cortex \shortcite{tompary_2017}}

\paragraph{Experimental goal.}
We tested whether offline consolidation increases the similarity of neural representations for episodes that share a common scene context. As in \shortciteA{tompary_2017}, overlap was evaluated before and after sleep-induced consolidation.

\paragraph{Episode set and latent structure.}
Each episode contained one object latent ($A_i$) and one scene latent ($B_j$). We used 128 distinct objects and 4 scenes. Objects were assigned evenly to scenes (32 objects per scene), yielding a deterministic training sequence in which each object was presented with its assigned scene.

\paragraph{Encoding protocol.}
Networks were initialized from scratch and trained for one wake day. Each object--scene episode was presented once, with fixed duration (5 timesteps per episode). Trial-level variability was introduced by swap noise that preserves global sparsity (equal number of active-to-inactive and inactive-to-active flips). This encoding regime produces one noisy realization per episode while keeping category-level structure intact.

\paragraph{Pre- and post-consolidation probing.}
Representational overlap was probed with partial object cues: object information was provided, whereas scene input channels were set to zero. For each cue, MTL cue was pattern completed, fixing the neural activity corresponding to the cue. The completed MTL state was then propagated through the feed-forward/feedback pathway to obtain the final evoked patterns in both MTL and CTX. This probe was performed immediately before sleep (\textit{pre}) and repeated after consolidation (\textit{post}).

\paragraph{Consolidation manipulation.}
Consolidation consisted of one sleep cycle including both replay streams (episodic and semantic replay). Semantic replay was enabled from the beginning of training (i.e., no initial phase with episodic-only replay). Replay durations followed the Figure experiments configuration listed in Table~\ref{tab:overlap_reproduction_parameters}.

\paragraph{Overlap metric and summary statistics.}
For each region (MTL, CTX), we computed a pairwise Pearson correlation matrix across all object-evoked probe patterns. For object $i$, we defined overlap as:
\begin{equation}
\mathrm{OverlapIndex}_i
=
\overline{r}_{i,\mathrm{same\ scene}}
-
\overline{r}_{i,\mathrm{different\ scene}},
\end{equation}
where $\overline{r}_{i,\mathrm{same\ scene}}$ is the mean correlation between object $i$ and other objects assigned to the same scene, and $\overline{r}_{i,\mathrm{different\ scene}}$ is the corresponding mean across objects assigned to different scenes. Positive values indicate stronger scene-dependent clustering. We summarized the distribution of per-object overlap indices by phase and region (mean, standard deviation, and standard error).
\begin{table}[ht]
\centering
\caption{Parameters used to reproduce consolidation-induced overlap.}
\label{tab:overlap_reproduction_parameters}
\begin{tabular}{p{0.58\linewidth}p{0.30\linewidth}}
\toprule
\textbf{Parameter} & \textbf{Value} \\
\midrule
Number of object concepts ($|A|$) & 128 \\
Number of scene concepts ($|B|$) & 4 \\
Object-to-scene assignment & Balanced (32 objects per scene) \\
Concept code size (object, scene) & 10 active units, 10 active units \\
Sensory/MTL-sensory size & Auto-allocated (1280 object + 40 scene units) \\
Wake training days & 1 \\
Presentations per object per day & 1 \\
Episode duration during encoding & 5 timesteps (fixed) \\
Training noise & 4 swap operations per timestep \\
Cue type for overlap probe & Object-only cue (scene channels silenced) \\
Cue noise at probe & 4 swap operations \\
MTL retrieval during probe & Pattern completion with cue-constrained sensory component \\
Number of consolidation cycles & 1 \\
Replay regime during consolidation & Episodic + semantic replay active \\
Sleep replay duration (episodic stream) & 500 timesteps \\
Sleep replay duration (semantic stream) & 50 timesteps \\
Random seed & 1 \\
\bottomrule
\end{tabular}
\end{table}

\subsubsection*{Neocortical synaptic engrams for remote contextual memories \shortcite{lee_2023}}

\paragraph{Experimental goal and circuit mapping.}
To reproduce the synaptic-engram findings of contextual fear conditioning (\shortciteA{lee_2023}), we mapped cortical and subcortical populations in the model to PFC-like (CTX) and BLA-like (MTL-semantic) representations, respectively. We then quantified how consolidation and extinction reshape synaptic interactions among tagged neuronal subpopulations.

\paragraph{Task design and latent structure.}
Episodes were defined over two latent attributes with 5 concepts each. We interpreted one valence concept ($B_1$) as fear (US) and one context concept ($A_1$) as conditioned context (CS). During conditioning, fear was persistently co-activated with variable context components, establishing a consolidated fear-associated trace. A subsequent recall phase probed reactivation under CS/US-matched input.

\paragraph{Conditioning and recall protocol.}
Conditioning consisted of one wake day followed by sleep consolidation. Recall was then run as a dedicated wake day, again followed by sleep. To match the reference protocol, one sensory channel was randomized within a restricted temporal window during recall, preserving the overall latent identity while introducing within-trial variability. An additional recall pass was used for recall-specific population tagging and downstream synaptic analyses.

\paragraph{Engram and recall-only population definitions.}
Engram populations were defined from mean wake activity during encoding-related windows using the model's sparse activation rule. In CTX, \emph{engram} cells were those active in the conditioning window; in MTL-semantic, engram cells were extracted from the first recall window. We also defined CTX \emph{recall-only} cells as units active during the later recall probe but not active during conditioning. Synaptic-strength distributions were then measured for: (i) CTX engram$\rightarrow$engram vs non-engram$\rightarrow$engram projections, (ii) MTL-semantic engram vs non-engram projections onto CTX engram cells, and (iii) CTX recall-only projections onto CTX engram cells.

\paragraph{Extinction and behavioural proxy.}
Extinction was modeled as repeated CS-context exposure without fear-US latent activation across many days, each day followed by sleep. A fear-expression proxy was computed after each extinction day as the recurrent CTX input onto a pre-identified fear-cell subset, evoked from an early recall-derived cortical state. This yielded a day-by-day trajectory of fear-related drive during extinction.

\paragraph{Post-extinction synaptic reorganization and savings analysis.}
After extinction, we re-estimated CTX engram-related synaptic distributions (engram$\rightarrow$engram and non-engram$\rightarrow$engram) to quantify extinction-associated weakening of engram-specific structure. To evaluate memory savings, we compared recall stability in the extinguished network against a conditioning-only control using pattern completion from matched probe states, with cosine similarity as the stability measure.
\subsection*{Semantic facilitation of recall}
To test whether semantic structure in the episodic memory circuit improves recall (Fig. \ref{fig:behavioural}c), we compared retrieval in networks with intact semantic support against networks in which the semantic pathway was disrupted or scrambled. The logic of this manipulation was to isolate the contribution of semantically interpretable input in episodic recall \shortcite{lin_2021}.
\newline\newline
We compared an \emph{intact} pretrained network against a \emph{scrambled} version in which semantic-to-episodic guidance was disrupted. Test input consisted of a single day of length $T_\mathrm{day}=25$ timesteps, corresponding to 5 episodes of fixed duration 5. All simulations in this panel were run with a single random seed.
\newline\newline
In each trial, the network was first run on this day-long input without sleep, and the resulting MTL activity patterns were stored as target episode representations. Recall was then probed separately for each stored episode by corrupting its MTL-sensory component with cue corruption $N_\mathrm{swaps}=5$, initializing the remaining MTL units to zero, and iterating recurrent MTL dynamics for 10 steps. Recall performance was quantified as the cosine similarity between the recovered full MTL state and the original stored episode. This procedure was repeated across 10 trials, where each trial consisted of a newly generated day-long input sequence followed by recall of all stored episodes.
\newline\newline
To construct the scrambled condition, we permuted the semantic-to-episodic weight matrix $W^{\mathrm{CTX}\leftarrow \mathrm{MTL\text{-}semantic}}$ independently across rows, thereby preserving the overall weight distribution while destroying the alignment between semantic structure and episodic assemblies. We additionally froze these scrambled synapses and suppressed semantic recruitment bias. Thus, the intact and scrambled networks were matched in overall architecture and synaptic magnitude, but differed in whether semantic input remained coherently aligned with the learned episodic code. Improved recall in the intact condition therefore reflects the contribution of structured semantic support to MTL pattern completion.
\subsection*{Blocked versus interleaved training}
To compare blocked and interleaved training schedules (Fig. \ref{fig:behavioural}d), we trained separate networks on the same set of simple concepts while varying only the temporal ordering of presentations. In both conditions, sensory input was corrupted with $N_\mathrm{swaps}=4$. In the blocked condition, training was organized into 10 separate runs of length $T_\mathrm{day}=50$ timesteps each, with one concept presented per run and fixed episode duration 1, such that concepts were encountered in temporally grouped blocks. In the interleaved condition, the network was instead trained for 100 days of the same length ($T_\mathrm{day}=50$) and duration statistics, with all 25 latent combinations sampled uniformly. All simulations in this panel were run with a single random seed.
\newline\newline
After training, replayed MTL-sensory patterns were collected from Sleep A. For the blocked condition, replay was sampled separately after each concept-specific training run; for the interleaved condition, replay was sampled after the full mixed training regime. We then constructed simple concept prototypes by averaging awake MTL-sensory activity within each concept condition and binarizing the resulting pattern with the MTL-sensory activation rule.
\newline\newline
For each replay event, we computed its cosine overlap with the full library of simple concept prototypes and retained the maximal overlap. Replay quality was quantified as the distribution of these maximal overlaps across replay events. Thus, this analysis isolates the effect of temporal training structure on the prototype-likeness of replay while holding the sensory corruption level fixed.
\subsection*{Anterograde lesion of MTL for semantic recall}
To model anterograde amnesia (Fig. \ref{fig:behavioural}f), we started from a pretrained network with pre-existing higher-order selectivity and then introduced a single new paired association $(A_i,B_i)$. Baseline input statistics were defined over one day of length $T_\mathrm{day}=80$ with sensory corruption $N_\mathrm{swaps}=2$, while the new association was presented in a dedicated association day of length 10. All simulations in this panel were repeated across 20 seeds.
\newline\newline
Before learning the new pair, the cortical episode subregion corresponding to previously absent diagonal pairs was reset to an immature state, so that the new association required fresh cortical recruitment. Two copies of the network were then exposed to the same single-day $(A_i,B_i)$ association. In the sham condition, the network was allowed to undergo normal sleep-dependent consolidation. In the anterograde condition, MTL recurrent consolidation machinery was lesioned immediately after the awake association day, by setting recurrent MTL connectivity to zero ($W^{\mathrm{MTL}\leftarrow \mathrm{MTL}}$, $W^{\mathrm{MTL\text{-}sensory}\leftarrow \mathrm{MTL\text{-}sensory}}$, and $W^{\mathrm{MTL\text{-}semantic}\leftarrow \mathrm{MTL\text{-}semantic}}$) and disabling Sleep A replay, before sleep was invoked.
\newline\newline
After this manipulation, recall was tested by cueing the learned $A_i$ component and measuring cortical recruitment of the target $B_i$ representation after 10 iterations of CTX pattern completion. The main readout was the mean activity of the target cortical $B_i$ assembly after the $A_i$ cue. To characterize robustness, we repeated this procedure while varying cue corruption of the $B$ component over $N_\mathrm{swaps}^B\in\{0,\dots,10\}$. This analysis therefore isolates the effect of disrupting post-encoding MTL replay and recurrent completion on the consolidation of a newly acquired association.
\subsection*{Retrograde lesion of MTL for semantic recall}
To model retrograde semantic amnesia (Fig. \ref{fig:behavioural}g), we started from a pretrained network with pre-existing higher-order selectivity and introduced a single new paired association $(A_i,B_i)$. Baseline input statistics were defined over one day of length $T_\mathrm{day}=80$ with sensory corruption $N_\mathrm{swaps}=2$. The new association was then presented for one day of length 25, and both sham and retrograde networks were allowed to undergo sleep-dependent consolidation before the lesion was applied. Simulations in the main comparison were repeated across 20 seeds.
\newline\newline
After this post-learning sleep phase, the retrograde condition was generated by lesioning the semantic pathway while leaving the rest of the network unchanged. Concretely, we removed semantic participation by setting the semantic lesion flag, zeroing semantic-to-episodic connectivity, removing semantic contributions to CTX input through the MTL$\rightarrow$CTX pathway, and eliminating recurrent semantic connectivity. The sham network was left intact.
\newline\newline
Recall was then probed using a partial sensory cue containing only the $A_i$ component of the learned pair. This cue first recruited the corresponding cortical context representation, after which additional cortical corruption was introduced by swapping $N_\mathrm{swaps}^{\mathrm{CTX}}$ units within CTX subregions. In the sham condition, the corrupted cortical state could then reactivate MTL-semantic and return semantic support to cortex; in the retrograde condition, this route was unavailable. Final cortical completion was obtained by 10 iterations of CTX pattern completion, and recall performance was quantified as the mean activity of the target $B_i$ cortical assembly after the $A_i$ cue.
\newline\newline
We repeated the analysis over a noise sweep with $N_\mathrm{swaps}^{\mathrm{CTX}}\in\{10,\dots,39\}$, using 10 seeds per noise level. This analysis isolates the contribution of semantic feedback to the recall of previously learned associations under degraded cortical cues.
\begin{table}[ht]
\centering
\caption{Parameters used for synaptic-engram reproduction.}
\label{tab:synaptic_engrams_parameters}
\begin{tabular}{p{0.58\linewidth}p{0.30\linewidth}}
\toprule
\textbf{Parameter} & \textbf{Value} \\
\midrule
Latent attributes and concepts & 2 attributes, 5 concepts each \\
Concept code size (context, valence) & 10 active units, 10 active units \\
Conditioning schedule & 1 day, 100 timesteps, fixed episode duration 5 \\
Conditioning noise & 4 swap operations per timestep \\
Conditioning latent statistics & Persistent fear-US ($B_1$) with variable context components \\
Recall schedule & 1 day, CS/US-matched probe distribution \\
Recall variability control & Restricted temporal channel randomization (protocol match) \\
Extinction schedule & 1000 days, 20 timesteps/day, fixed episode duration 5 \\
Extinction noise & 4 swap operations per timestep \\
Extinction latent statistics & CS-context presentations without fear-US latent \\
Fear-probe index within recall trajectory & 25\% of recall window \\
Engram tagging basis & Mean wake activity with sparse activation criterion \\
Synaptic comparisons & CTX E$\rightarrow$E vs NE$\rightarrow$E; MTL-semantic E/NE$\rightarrow$CTX E; CTX recall-only$\rightarrow$CTX E \\
Savings metric & Cosine similarity after CTX pattern completion \\
Number of savings seeds & 50 \\
Recording rate for wake indexing & 1 (full-rate recording) \\
Random seed (reported notebook run) & 0 \\
\bottomrule
\end{tabular}
\end{table}

\section*{Acknowledgments}
We would like to thank Rachel Swanson and Colleen J. Gillon, as well as the rest of the members of the Clopath Lab, for discussions and  comments on earlier versions of the manuscript. We gratefully acknowledge the financial support of the Wellcome Trust, the Simons Foundation, the Engineering and Physical Sciences Research Council (EPSRC), and the European Research Council (ERC).
\section*{Declaration of Interests}
The authors declare no competing interests.
\section*{Code Availability}
All code supporting the findings of this study is available at 
\url{https://github.com/albesagonzalez/sensory-semantic-episodes}
\section*{Declaration of generative AI and AI-assisted technologies in the manuscript preparation process}
During the preparation of this work the authors used ChatGPT in order to assist with writing at the sentence level, and coding at the snippet level. After using this tool/service, the authors reviewed and edited the content as needed and take full responsibility for the content of the published article.
\bibliography{Cites.bib}
\clearpage
\section*{Supplementary Material}
% redefine figure numbering for supplements
\renewcommand{\thefigure}{S\arabic{figure}}
\setcounter{figure}{0} % reset counter so it starts at S1

\begin{figure}[h!]
    \centering
    \includegraphics[width=0.95\linewidth]{Figures/Figure_S1.png}
    \caption{\textbf{Episode generation protocol}
\textbf{a1}: Episodes are generated by first sampling latent concepts according to the predefined episode probability distribution $P(A_i,B_j)$. Each sampled episode corresponds to a combination of semantic factors that define the underlying event.
\textbf{a2}: The sampled latent concepts activate their associated sensory features, producing a de-noised input pattern. In the example shown, the episode contains two concepts (illustrated as pepperoni and cheese), each represented by a dedicated subset of active neurons.
\textbf{a3}: Sensory variability is introduced through swap noise. A total of $N_{\mathrm{swaps}}$ active neurons are randomly deactivated and $N_{\mathrm{swaps}}$ inactive neurons are randomly activated, generating a noisy sensory realization of the same latent episode. This procedure preserves the underlying semantic content while controlling the input signal-to-noise ratio (SNR).}
    \label{fig:input_generation}
\end{figure}
\begin{figure}[h!]
    \centering
    \includegraphics[width=0.95\linewidth]{Figures/Figure_S2.png}
    \caption{Connectivity matrices across the circuit (after learning)}
    \label{fig:circuit}
\end{figure}
\begin{figure}[h!]
    \centering
    \includegraphics[width=0.95\linewidth]{Figures/Figure_S3.png}
    \caption{\textbf{Comparison of recurrent cortical connectivity (circuit model) versus semantic structure (cognitive model). \textbf{a}}: Comparison after simple concepts vocabulary has been formed. \textbf{b}: Comparison after initial vocabulary has been extended to include complex concepts.}
    \label{fig:recurrent_semantics}
\end{figure}
\begin{figure}[h!]
    \centering
    \includegraphics[width=0.95\linewidth]{Figures/Figure_S4.png}
    \caption{\textbf{a}: Schematics of cortical dictionary update after MTL-semantic is developed. Initially, CTX$\leftarrow$MTL-sensory drives CTX activity during episodic replay. Because there is structure in MTL-semantic, this slowly updates CTX$\leftarrow$MTL-semantic connectivity. At learning convergence, this allowing semantic indexation from MTL to CTX during episodic replay.  \textbf{b}: Example MTL activity during episodic replay when reaplys is driven by full recurrent connectivity (MTL-sensory and MTL-semantic are cuopled via reciprocal connections too). \textrm{c}: CTX$\leftarrow$ MTL connectivity after learning with full MTL episodic replay.}
    \label{fig:phase_b_supp}
\end{figure}
\begin{figure}[h!]
    \centering
    \includegraphics[width=0.95\linewidth]{Figures/Figure_S5.png}
    \caption{\textbf{Generalization of replayed patterns and compositional generalization.}
\textbf{a}: Relationship between cortical test-time generalization and replay-to-awake generalization. Each point corresponds to an independent network realization, including different types of networks. Networks in which replayed representations more closely match their corresponding awake representations exhibit higher test-time generalization performance ($r=0.86$, $R^2=0.75$).
\textbf{b}: Relationship between cortical test-time generalization and replay-to-replay generalization. Networks with more consistent replay events corresponding to the same latent episode show improved generalization performance ($r=0.95$, $R^2=0.91$). The stronger correlation compared to panel a suggests that replay reliability across replay events is a key determinant of successful compositional consolidation.}
    \label{fig:normative_supp}
\end{figure}
\begin{figure}[h!]
    \centering
    \includegraphics[width=0.95\linewidth]{Figures/Figure_S6.png}
    \caption{Schematics of an experimental paradigm designed to test the predictions of the model. Compositional visual-odor cues appear following a specific probability in a treadmill, and are associated with either a reward or a negative reinforcement.}
    \label{fig:predictions_supp}
\end{figure}
\end{document}